% !LW recipe=latexmk (xelatex)
% Tesina de grado -- Licenciatura en Ciencias de la Computación
% Autor: Facundo Tomás Gaviola
%
% Basado en la plantilla UMU Thesis v1.1.0 (Enrique Pérez Navarro), adaptación de ipleiria-thesis
% (José Areia). La plantilla original, sin modificar, está en Tesis_Universidad_de_Murcia__UMU_/.
%
% Compilación: XeLaTeX con -shell-escape (minted) y biber.
%     latexmk -xelatex -shell-escape tesina.tex
%
% Estado: ESTRUCTURA. Cada capítulo contiene sus secciones con un recuadro "Pendiente" (\pendiente) que
% describe el contenido previsto. Las fuentes de cada capítulo y la extensión estimada figuran como
% comentarios (% Fuentes: / % Extensión estimada:). Los recuadros se eliminan a medida que se redacta.

%%% Opciones del documento %%%
\documentclass[
    language=spanish,
    docstage=working,
    media=screen,
    linkcolor=black,
    chapterstyle=fancy,
]{UMUthesis}

\DocumentVersion{0.1.0 (estructura)}

%%% Marcadores de contenido pendiente (\pendiente, \listadependientes) %%%
\usepackage{Configurations/12-Pendientes}

%%% Estilos compartidos por los diagramas de Figures/Capitulos %%%
\usepackage{Configurations/13-Figuras}

%%% Metadatos %%%
% Autor
\FirstAuthor{Facundo Tomás Gaviola}

% Director
\Supervisor{Pablo Vidal}
\SupervisorMail{pablo.vidal@ingenieria.uncuyo.edu.ar}

% Título en inglés
\Title{Detection, diagnosis and automatic report generation for skin lesions using vision language models}

% Título en español
\TitleEsp{Detección, diagnóstico y generación automática de reportes para lesiones de piel utilizando modelos de lenguajes de visión}

% Universidad, carrera y año
\University{Universidad Nacional de Cuyo}
\Degree{Licenciatura en Ciencias de la Computación}

% Lugar y fecha
\Date{Mendoza, \DTMmonthname{\month} \number\year}


%%% Glosario y acrónimos %%%
\makeglossaries
\loadglsentries{Matter/04-Glossary}
\loadglsentries[\acronymtype]{Matter/05-Acronyms}
% Lista todos los acrónimos definidos aunque todavía no se usen con \gls en el texto.
% Requiere que latexmk ejecute makeglossaries (configurado en .latexmkrc).
\glsaddall[types={\acronymtype}]

\begin{document}

%%% Portada %%%
% ----------------------------------------
% Carátula: fondo blanco con el isologotipo de la UNCuyo y la Facultad de Ingeniería. Contenido según el reglamento de
% la Tesina Final de Carrera (Res. CD 107/17, §3.2): título, alumno, tutor y fecha de presentación.
% El isologotipo es la versión oficial en color descargada de
% https://cicunc.uncuyo.edu.ar/imagen-institucional (marcas-facu-ingenieria.pdf), recortada a la marca combinada.
% ----------------------------------------
\newgeometry{
  top=2cm, bottom=2cm,
  left=2.5cm, right=2cm,
  bindingoffset=0cm
}
\begin{titlepage}
  \begingroup
    \plexsans % Aplica IBM Plex Sans solo localmente
    \color{frontpagedark}

    \noindent\includegraphics[width=10cm]{Figures/Theme/isologo-uncuyo-fing.pdf}\par
    \vspace{0.2cm}
    {\noindent\color{uncuyoazul}\rule{\textwidth}{0.8pt}\par}

    \vspace*{\fill} % empuje vertical

    % Selección de portada en inglés vs español
    \ifthenelse{\equal{\LanguageOption}{english}}{%
      {\noindent\bfseries\fontsize{12}{14}\selectfont FINAL DEGREE THESIS\par}
      \vspace{0.15cm}
      {\noindent\fontsize{11}{14}\selectfont Bachelor's Degree in Computer Science\par}
      \vspace{0.5cm}
      {\noindent\raggedright\itshape\color{uncuyoazul}\fontsize{24}{29}\selectfont \thetitle\par}
      \vspace{2cm}
      \def\EtiquetaAutor{AUTHOR:}\def\EtiquetaDirector{SUPERVISOR:}
    }{%
      {\noindent\bfseries\fontsize{12}{14}\selectfont TESINA FINAL DE CARRERA\par}
      \vspace{0.15cm}
      {\noindent\fontsize{11}{14}\selectfont \GetDegree\par}
      \vspace{0.5cm}
      {\noindent\raggedright\itshape\color{uncuyoazul}\fontsize{24}{29}\selectfont \titleesp\par}
      \vspace{2cm}
      \def\EtiquetaAutor{AUTOR:}\def\EtiquetaDirector{TUTOR:}
    }

    \noindent
    \begin{tabbing}
      \hspace{3cm}\=\kill
      \textbf{\fontsize{12}{10}\selectfont \EtiquetaAutor}\> \GetFirstAuthor \\[0.3cm]
      \textbf{\fontsize{12}{10}\selectfont \EtiquetaDirector}\> \GetSupervisorTitle~\GetSupervisor
    \end{tabbing}

    \vspace{1cm}
    {\noindent\hfill\fontsize{12}{16}\selectfont Mendoza, \DTMmonthname{\month} de \number\year\par}

    \vspace{1cm}
  \endgroup
\end{titlepage}
\restoregeometry
\MediaOptionLogicBlank

% Portadilla interior de la plantilla UMU. Las tesinas de la Facultad de Ingeniería llevan una sola carátula.
% % ----------------------------------------
% Definición del fondo 
% ----------------------------------------
  \newcommand\BackgroundPicCoverWhite{%
    \put(0,0){%
      \parbox[b][\paperheight]{\paperwidth}{%
        \vfill
        \centering
        \includegraphics[width=\paperwidth,height=\paperheight,keepaspectratio]{Figures/Theme/BG_UMU_White.pdf}%
        \vfill
      }%
    }%
  }%


\AddToShipoutPictureBG*{\BackgroundPicCoverWhite}

% ----------------------------------------
% Portada con cambio según LanguageOption
% ----------------------------------------
\newgeometry{
  top=1.5cm, bottom=1.77cm,
  left=1.5cm, right=1.5cm,
  bindingoffset=0cm
}
\begin{titlepage}
  \AddToShipoutPictureBG*{\BackgroundPicCoverWhite}
  \begingroup
    \plexsans % Aplica IBM Plex Sans solo localmente
    % Color de texto según opción de portada
    \color{black}

\vspace*{\fill}  % empuje vertical

    % Selección de portada en inglés vs español
\ifthenelse{\equal{\LanguageOption}{english}}{%
      % ======= PORTADA EN INGLÉS =======
      {\noindent\bfseries\fontsize{12}{14}\selectfont DOCTORAL THESIS\par}
      \vspace{0.5cm}
      {\noindent\raggedright\itshape\fontsize{22}{24}\selectfont \titleesp\par}
      \vspace{2cm}

      \noindent
      \begin{minipage}[t]{0.40\textwidth}
        \raggedright
        {\itshape\fontsize{14}{20}\selectfont \thetitle\par}
      \end{minipage}%
      \begin{minipage}[t]{0.4\textwidth}
        \raggedright
        \begin{tabbing}
          \hspace{4cm}\=\kill
          \textbf{\fontsize{12}{10}\selectfont AUTHOR:}\> \GetFirstAuthor \\[0.3cm]
          \textbf{\fontsize{12}{10}\selectfont SUPERVISORS:}\> \GetSupervisor
        \end{tabbing}
      \end{minipage}%
      \begin{minipage}[t]{0.1\textwidth}
        \vspace{1.5cm}
        \raggedleft
        {\fontsize{13}{16}\selectfont \the\year \par}
      \end{minipage}
    }{%
      % ======= PORTADA EN ESPAÑOL =======
      {\noindent\bfseries\fontsize{12}{14}\selectfont TESINA DE GRADO — LICENCIATURA EN CIENCIAS DE LA COMPUTACIÓN\par}
      \vspace{0.5cm}
      {\noindent\raggedright\itshape\fontsize{22}{24}\selectfont \titleesp\par}
      \vspace{2cm}

      \noindent
      
      \begin{minipage}[t]{0.4\textwidth}
        \raggedright
        \begin{tabbing}
          \hspace{4cm}\=\kill
          \textbf{\fontsize{12}{10}\selectfont AUTOR:}\> \GetFirstAuthor \\[0.3cm]
          \textbf{\fontsize{12}{10}\selectfont DIRECTOR:}\> \GetSupervisor
        \end{tabbing}
      \end{minipage}%
      \begin{minipage}[t]{0.45\textwidth}
        \vspace{2cm}
        \raggedleft
        {\fontsize{13}{16}\selectfont \the\year \par}
      \end{minipage}
    }
   % Fin condicional LanguageOption

    \vspace{1cm}
  \endgroup
\end{titlepage}
\restoregeometry
\MediaOptionLogicBlank

%%% Declaración de autoría %%%
% La plantilla incluye la declaración de la UMU. Reemplazar por la que exija la UNCuyo (si corresponde)
% antes de reactivar esta línea.
% \pagenumbering{gobble} % Prevent page numbering.

\thispagestyle{plain} % Page style without header and footer

% El nombre de este fichero debe ser ASCII puro. Los acentos en nombres de
% fichero sobreviven en macOS (APFS ignora la normalizacion Unicode) pero se
% rompen en Linux/Overleaf, donde NFC y NFD son nombres distintos.
\includepdf[pages={1}]{Documents/DECLARACION_DE_AUTORIA_Y_ORIGINALIDAD.pdf}

\MediaOptionLogic
\pagenumbering{gobble}

%%% Numeración romana %%%
\pagenumbering{roman}

%%% Dedicatoria y agradecimientos %%%
\chapter*{} % Sin numeración
\thispagestyle{empty}

\begin{flushright}
    \textit{Dedicatoria (a completar).}
\end{flushright}

\chapter*{Agradecimientos}

\pendiente{A completar al final de la redacción}


%%% Resumen / Abstract %%%
\thispagestyle{plain}
\chapter*{Resumen}
% Extensión: unas 340 palabras. Cifras tomadas de los Caps. 4 y 5 (tab:res_supervisados, tab:res_pruebas_vlm,
% sec:res_melanoma, sec:res_fuga, sec:res_hiba, tab:res_gen_cascada, tab:res_gen_pipeline, sec:res_gen_concordancia,
% tab:res_costo_generadores). El abstract repite la misma estructura y las mismas cifras.

El diagnóstico de una lesión de piel a partir de una imagen dermatoscópica exige, además de clasificarla, redactar un
reporte que registre lo observado y la conclusión. Los modelos visión--lenguaje (VLM) prometen resolver las dos tareas
con un mismo modelo, pero no está establecido si clasifican las lesiones tan bien como los clasificadores entrenados
solo con imágenes, ni si los reportes que generan afirman el diagnóstico correcto y describen la lesión de la imagen.
Este trabajo estudia ambas preguntas sobre \emph{HAM10000}, con una partición disjunta por lesión, y bajo la
restricción de usar una sola GPU de consumo.

En la clasificación se comparan cuatro clasificadores supervisados, dos VLM generativos de propósito general y un
modelo visión--lenguaje dermatológico, DermLIP. El mejor supervisado, Swin-Tiny, obtiene un F1 macro de
$0.733 \pm 0.026$ en la prueba. Los VLM generativos quedan cerca del azar, y DermLIP sin ajuste queda entre 0.15 y
0.18 de sensibilidad macro por debajo de los mejores supervisados, aunque con una cabeza entrenada sobre su
representación se acerca a ellos. En melanoma, un umbral fijado en validación lleva la detección de Swin-Tiny de tres a
cuatro de cada cinco casos. La evaluación muestra, además, que una partición por imagen sobreestima la sensibilidad
macro en unas 0.07, y que en otra población, \emph{HIBA}, el clasificador pierde cerca del 40\,\% de ella.

Para la generación se propone un \emph{pipeline} que entrega al generador el diagnóstico de Swin-Tiny y decide, según
la concordancia con un \emph{k}-vecinos sobre la representación de DermLIP, si afirmarlo, presentar dos candidatos o
derivar el caso a un especialista. El reporte transcribe el diagnóstico recibido, de modo que su consistencia
diagnóstica, 0.849, coincide con la exactitud del clasificador, y el ruteo la eleva a 0.870. Entre los casos que el
\emph{pipeline} afirma, el clasificador acierta el 91.5\,\%, aunque los errores que comparten los dos predictores se
siguen afirmando. El generador, Mini-VLM, ocupa como máximo 8.6~GB y escribe un reporte en unos 16 segundos.

El trabajo establece la fidelidad diagnóstica de los reportes con métricas automáticas, pero no su validez clínica,
que requiere la evaluación de especialistas, y sus resultados se limitan a la dermatoscopía.

\keywordses{lesiones de piel, dermatoscopía, modelos visión--lenguaje, clasificación diagnóstica, detección
de melanoma, generación de reportes clínicos, incertidumbre}

\pdfbookmark[1]{Abstract}{abstract}
\chapter*{Abstract}

Diagnosing a skin lesion from a dermoscopic image requires not only classifying it but also writing a report that
records the findings and the conclusion. Vision--language models (VLMs) promise to solve both tasks with a single
model, but it is not established whether they classify lesions as well as classifiers trained on images alone, nor
whether the reports they generate state the correct diagnosis and describe the lesion in the image. This work studies
both questions on \emph{HAM10000}, with a lesion-disjoint split, under the constraint of a single consumer GPU.

For classification, four supervised classifiers, two general-purpose generative VLMs and a dermatology
vision--language model, DermLIP, are compared. The best supervised model, Swin-Tiny, reaches a macro F1 of
$0.733 \pm 0.026$ on the test set. The generative VLMs stay close to chance, and zero-shot DermLIP falls 0.15 to 0.18
of macro sensitivity below the best supervised models, although with a head trained on its representation it comes
close to them. For melanoma, a threshold set on validation raises Swin-Tiny's detection from three to four out of five
cases. The evaluation also shows that an image-level split overestimates macro sensitivity by about 0.07, and that on
another population, \emph{HIBA}, the classifier loses about 40\,\% of it.

For report generation, a pipeline is proposed that passes the Swin-Tiny diagnosis to the generator and decides,
according to its agreement with a \emph{k}-nearest-neighbours vote on the DermLIP representation, whether to assert
it, present two candidates or refer the case to a specialist. The report transcribes the diagnosis it receives, so
its diagnostic consistency, 0.849, matches the classifier's accuracy, and routing raises it to 0.870. Among the cases
the pipeline asserts, the classifier is correct in 91.5\,\%, although the errors shared by both predictors are still
asserted. The generator, Mini-VLM, uses at most 8.6~GB and writes a report in about 16 seconds.

This work establishes the diagnostic fidelity of the reports with automatic metrics, but not their clinical validity,
which requires evaluation by specialists, and its results are limited to dermoscopy.

\keywordsen{skin lesions, dermoscopy, vision--language models, diagnostic classification, melanoma
detection, clinical report generation, uncertainty}


%%% Índices %%%
\bookmarktocentry\tableofcontents
\listoffigures
\listoftables
\listadependientes % Solo con docstage=working.

%%% Acrónimos %%%
\acronymtoc\printglossary[type=\acronymtype,title={Acrónimos},nonumberlist]

%%% Numeración arábiga %%%
\cleardoublepage
\pagenumbering{arabic}

%%% Capítulos %%%
\chapter[Introducción]{Introducción}
\label{cp:introduccion}
% Extensión estimada: sin límite; la ficha orienta 6--7 páginas (ver fichas/01-introduccion.md y guía §5).
% Fuentes: Plan de trabajo (Introducción y objetivos), Propuesta (secciones "Introducción y Propósito" y
% "Marco de la Problemática"), PLAN_TFG_GAVIOLA_FACUNDO.pdf y la revisión del tutor del 28/09/2026
% (tareas_revision_cap1-2.md, I1--I8). Sin hipótesis (D1): solo preguntas de investigación y objetivos.

Este capítulo presenta el problema que motiva el trabajo, las preguntas de investigación y los objetivos, y
describe cómo se organiza el documento.

\section{Contexto y motivación}
\label{sec:intro_contexto}

El cáncer de piel es uno de los tipos de cáncer más frecuentes del mundo. La última estimación de GLOBOCAN,
publicada en 2026, atribuye al año 2024 1\,138\,314 casos nuevos de cáncer de piel no melanoma y 337\,969 de
melanoma~\cite{sung2026global}. El melanoma es menos de una cuarta parte de los diagnósticos, pero causa casi
tantas muertes como el resto de los cánceres de piel juntos, 58\,398 frente a 58\,721.

El desenlace del melanoma depende del momento en que se lo diagnostica. La supervivencia relativa a cinco años
es de 100.0\,\% cuando se detecta localizado, baja a 76.0\,\% cuando alcanzó los ganglios regionales y a
34.0\,\% cuando hay metástasis a distancia~\cite{nci2026melanoma}. Por eso importan las herramientas de apoyo al
diagnóstico que permitan detectarlo a tiempo.

La detección temprana enfrenta varias limitaciones. Los especialistas son escasos: un relevamiento en 158
países estimó una media de 2.66 dermatólogos por cada 100\,000 habitantes, con un rango de 0.37 en los países de
ingreso bajo a 5.05 en los de ingreso alto~\cite{freeman2026global}. Donde faltan, la primera valoración queda en
manos de clínicos que necesitan la dermatoscopía, cuyo beneficio depende de una formación que no todos
tienen~\cite{dinnes2018dermoscopy}. Aun con esa formación, la lectura de una misma imagen varía entre
observadores según su nivel de experiencia~\cite{tschandl2019comparison}. A esto se suma la documentación, que
consume una parte grande del tiempo médico~\cite{arndt2017tethered}.

Esas limitaciones retrasan la atención en dos puntos, la interpretación de la lesión y la redacción de su
reporte, y de ahí surgen dos tareas computacionales: asignar una categoría a la lesión y redactar su reporte.

Para la primera se aplicaron métodos de aprendizaje profundo, como las \glspl{cnn} y los \glspl{vit} entrenados
de forma supervisada sobre imágenes etiquetadas~\cite{dosovitskiy2021image,liu2021swin}. Ya en 2017 una \gls{cnn}
alcanzó un desempeño comparable al de dermatólogos certificados~\cite{esteva2017dermatologistlevel}.
Sin embargo, las revisiones sistemáticas registran exactitudes altas en condiciones
controladas~\cite{magalhaes2024systematic,kreouzi2024learning}, y queda abierto si ese desempeño se conserva en
atención primaria~\cite{jones2022intelligence}. Además, un clasificador produce una categoría y no el texto que
la describe, de modo que la documentación queda sin asistencia.

Los \glspl{vlm} responden a esa limitación porque integran información visual y textual. Se entrenan con pares
imagen--descripción que ubican ambas modalidades en un espacio común, lo que permite clasificar una imagen sin
fijar de antemano las categorías~\cite{radford2021learning}. Los modelos generativos agregan un modelo de
lenguaje que produce texto libre condicionado a la imagen~\cite{liu2023visual}, y su adaptación al dominio
biomédico mostró que el procedimiento se traslada a imágenes médicas~\cite{li2023llavamed}. Pueden ser generales
o estar especializados en un dominio.

Quedan abiertos dos temas. El primero es la generalización. El desempeño de los sistemas de \gls{ia}
dermatológica cae en los tonos de piel menos representados en los datos de entrenamiento, que dominan los
conjuntos públicos~\cite{daneshjou2022disparities,groh2021neural}, y la población argentina tampoco se puede
suponer homogénea, porque el Registro Argentino de Melanoma Cutáneo documentó diferencias regionales en las
características de la enfermedad~\cite{loria2020cutaneous}. El segundo son los reportes diagnósticamente fieles,
que describen la lesión que la imagen muestra y afirman el diagnóstico correcto sin sostener lo que la imagen no
respalda. En este trabajo se denomina fidelidad diagnóstica a la coincidencia entre el diagnóstico que el
reporte afirma y el de referencia de la imagen, y calidad descriptiva a la adecuación del texto a la lesión que
la imagen muestra.

Las dos tareas están conectadas, porque el diagnóstico de la clasificación es lo que el texto debe afirmar. Un
texto fluido puede sostener un diagnóstico equivocado, y un diagnóstico acertado puede acompañarse de una
descripción que no corresponde a la lesión, por lo que no basta resolverlas por separado. El problema de
investigación es entonces el estudio conjunto de la clasificación de lesiones y de la generación de reportes con
\glspl{vlm}, con recursos de hardware limitados.

\section{Preguntas de investigación}
\label{sec:intro_preguntas}
El problema se traduce en dos preguntas, una por tarea, porque sus salidas, una categoría y un texto, se evalúan
con métricas que no se traducen entre sí.

\begin{enumerate}[label=\textbf{PI\arabic*.}]
  \item ¿Cómo se compara el desempeño de \glspl{vlm} generales y especializados con el de clasificadores
  visuales supervisados en la clasificación de lesiones de piel y en la detección de melanoma?
  \item ¿Qué fidelidad diagnóstica y qué calidad descriptiva alcanzan, por separado, los reportes que genera un
  \gls{vlm} ajustado a partir de la imagen?
\end{enumerate}

\section{Objetivos}
\label{sec:intro_objetivos}
El objetivo general es diseñar, implementar y evaluar un enfoque de aprendizaje profundo y de \glspl{vlm},
adaptados a partir de modelos preentrenados, para clasificar lesiones cutáneas y generar reportes
dermatológicos a partir de imágenes, y analizar su desempeño diagnóstico, la fidelidad diagnóstica de los
reportes y su viabilidad computacional con recursos de hardware limitados.

Los objetivos específicos son:
\begin{enumerate}[label=\textbf{OE\arabic*.}]
  \item Diseñar, implementar y evaluar, adaptando modelos preentrenados, clasificadores de lesiones cutáneas
  basados en \glspl{cnn}, \glspl{vit} y \glspl{vlm} para identificar posibles melanomas, y comparar los
  \glspl{vlm} generales y especializados con los clasificadores supervisados.
  \item Diseñar, implementar y evaluar un \gls{vlm}, adaptado de un modelo preentrenado, que integre información
  visual y textual para producir un diagnóstico preliminar y el reporte de la lesión, con evaluación separada de
  la fidelidad diagnóstica y la calidad descriptiva.
  \item Evaluar los modelos desarrollados en distintas condiciones: la generalización a otra población, la
  calibración de la confianza y el tratamiento de la incertidumbre, y el costo computacional.
\end{enumerate}

La Tabla~\ref{tab:intro_objetivos_preguntas} relaciona cada objetivo específico con la pregunta que contribuye a
responder y con los capítulos donde se lo describe y se lo evalúa.

\begin{table}[htbp]
  \centering
  \caption{Relación entre los objetivos específicos, las preguntas de investigación y los capítulos.}
  \label{tab:intro_objetivos_preguntas}
  \small
  \begin{tabularx}{\textwidth}{@{}p{0.29\textwidth}p{0.10\textwidth}XX@{}}
    \toprule
    Objetivo específico & Pregunta & Se diseña y describe en & Se evalúa y responde en \\
    \midrule
    OE1. Clasificadores de lesiones y de melanoma (\gls{cnn}, \gls{vit}, \gls{vlm}) & PI1 &
      Capítulo~\ref{cp:metodos}, Materiales y métodos &
      Capítulo~\ref{cp:res_clasif}, Resultados de clasificación \\
    \filasep
    OE2. \gls{vlm} para generar reportes & PI2 &
      Capítulo~\ref{cp:metodos}, Materiales y métodos &
      Capítulo~\ref{cp:res_gen}, Resultados de generación \\
    \filasep
    OE3. Evaluación en distintas condiciones & PI1 y PI2 &
      Capítulo~\ref{cp:metodos} (protocolo de evaluación) &
      Capítulos~\ref{cp:res_clasif} y~\ref{cp:res_gen} (resultados) y Capítulo~\ref{cp:discusion} (discusión y
      limitaciones) \\
    \bottomrule
  \end{tabularx}
\end{table}

\section{Organización del documento}
\label{sec:intro_organizacion}
Después de esta introducción, el documento se organiza en seis capítulos:

\begin{description}
  \item[Capítulo~\ref{cp:marco}, Marco teórico.] Este capítulo presenta los conceptos necesarios para leer el
  resto y el estado del arte.
  \item[Capítulo~\ref{cp:metodos}, Materiales y métodos.] Este capítulo presenta los conjuntos de datos, los
  modelos, el \emph{pipeline} y el protocolo de evaluación.
  \item[Capítulo~\ref{cp:res_clasif}, Resultados de clasificación.] Este capítulo presenta la comparación de las
  estrategias de clasificación y la evaluación fuera del conjunto de entrenamiento.
  \item[Capítulo~\ref{cp:res_gen}, Resultados de generación.] Este capítulo presenta la evaluación de los
  reportes generados y el costo de cómputo de las configuraciones del \emph{pipeline}.
  \item[Capítulo~\ref{cp:discusion}, Discusión.] Este capítulo presenta la interpretación conjunta de los
  resultados, la respuesta a las preguntas de investigación y las limitaciones.
  \item[Capítulo~\ref{cp:conclusiones}, Conclusiones.] Este capítulo presenta los hallazgos principales, el
  balance de los objetivos y el trabajo futuro.
\end{description}

Los apéndices reúnen el material propio que no entra en el cuerpo, y el anexo, las fichas de los conjuntos de
datos y de los modelos externos que el trabajo utiliza.

\chapter[Marco teórico]{Marco teórico}
\label{cp:marco}
% Extensión estimada: a revisar. El reparto de la ficha daba 27--30 páginas, pero el bloque A cerró en 7
% contra las 10--11 previstas. Ver el encabezado de fichas/02-marco-teorico.md.
% Fuentes: Docs/marco_teorico/Documento_LLMs_VLLMs.tex (recortar NLP general, one-hot, RNN/LSTM a una
% mención breve), Docs/estado_del_arte/resumen_papers.tex (fichas de ~40 trabajos) y
% Docs/tabla_resultados_sota. Los fundamentos clínicos, de calibración y de estadística no están en esos
% documentos.
%
% Feedback del tutor (14/09/2026): el estado del arte pasa a ser la última sección de este capítulo, porque
% revisa y usa los conceptos introducidos antes. Los conjuntos de datos y el sesgo, que antes estaban en el
% estado del arte, son conceptos y por eso van aquí. Se agrega una subsección sobre generación de texto a
% partir de imágenes médicas, que faltaba pese a ser el objeto central del trabajo. La subsección "Uso de un
% VLM generativo como clasificador" se traslada a Materiales y métodos por tratarse del enfoque propio.
%
% Revisión 25/09/2026: la atención pasa a sección propia, antes de visión y lenguaje, porque la usan los
% dos (ficha, decisión 13). El capítulo queda solo con conceptos de computación: los fundamentos clínicos
% se reparten entre Materiales y métodos y las métricas de texto, y los conjuntos concretos pasan a
% Materiales y métodos; aquí quedan el sesgo y la generalización como conceptos (ficha, decisión 14).
% Se agregan los fundamentos del aprendizaje profundo como §2.1 y el capítulo sigue el recorrido del
% pipeline (ficha, «Propuesta del 25/09/2026», aprobada).

Este capítulo reúne los conceptos que usa el resto del trabajo, en el orden en que los atraviesa un caso
dentro del \emph{pipeline} (Figura~\ref{fig:marco_mapa}). Una imagen dermatoscópica llega a un
clasificador, que le asigna una categoría, y a un codificador visual, que recupera casos parecidos; con
esa información, un modelo de lenguaje redacta el reporte. La salida se juzga después con métricas de
clasificación y de texto, y queda por saber si lo medido se sostiene fuera de los datos de entrenamiento.
Los tres componentes son redes neuronales entrenadas con el mismo procedimiento, y casi todos se
construyen sobre la operación de atención; por eso el capítulo empieza por ahí. La última sección revisa
el estado del arte con esos conceptos ya disponibles.

% Figura: mapa del capítulo sobre el pipeline.
% Se usa en la apertura del Capítulo 2. Cada componente lleva la sección que lo explica. Es la versión
% simplificada del diagrama del pipeline de sec:met_pipeline: si uno cambia, revisar el otro.

\begin{figure}[htbp]
    \centering
    \begin{tikzpicture}[x=1mm, y=1mm]

        \node[figentrada, text width=26mm] (img) at (-4, 0) {Imagen\\dermatoscópica};

        \node[figtarea, text width=30mm] (clf) at (40, 13)
            {Clasificador\\{\footnotesize visión, §\ref{sec:marco_vision}}};
        \node[figtarea, text width=30mm] (knn) at (40, -13)
            {Codificador visual y kNN\\{\footnotesize §\ref{sec:marco_vlm}}};

        \node[figtarea, text width=36mm] (gen) at (86, 0)
            {Generador\\{\footnotesize lenguaje, §\ref{sec:marco_llm};\\visión--lenguaje, §\ref{sec:marco_vlm}}};

        \node[figsalida, text width=18mm] (rep) at (126, 0) {Reporte};

        \node[fignodo, text width=30mm, draw=figpurpura!70!black, fill=figpurpura!12] (eval)
            at (126, -24) {Evaluación, §\ref{sec:marco_evaluacion}\\{\footnotesize generalización,
            §\ref{sec:marco_generalizacion}}};

        \draw[figflecha] (img.east) -- ++(5, 0) |- (clf.west);
        \draw[figflecha] (img.east) -- ++(5, 0) |- (knn.west);
        \draw[figflecha] (clf.east) -- node[figetiqueta, above] {categoría} ++(12, 0) |- ([yshift=3mm]gen.west);
        \draw[figflecha] (knn.east) -- node[figetiqueta, below, align=center] {casos\\parecidos} ++(12, 0) |- ([yshift=-3mm]gen.west);
        \draw[figflecha] (gen.east) -- (rep.west);
        \draw[figflecha] (rep.south) -- (eval.north);

        \begin{scope}[on background layer]
            \node[figgrupo={todos los componentes: fundamentos, §\ref{sec:marco_fundamentos};
                  atención, §\ref{sec:marco_atencion}}, fit=(clf) (knn) (gen)] {};
        \end{scope}

    \end{tikzpicture}
    \caption[El capítulo recorrido sobre el \emph{pipeline}]{El capítulo recorrido sobre el
    \emph{pipeline}. Cada componente lleva la sección que presenta los conceptos que usa. El clasificador
    asigna una categoría a la imagen, un codificador visual recupera casos parecidos y el generador redacta
    el reporte con esa información; después la salida se evalúa y se examina si lo medido se sostiene
    fuera de los datos de entrenamiento. Los fundamentos del entrenamiento y la operación de atención son
    comunes a los tres componentes.}
    \label{fig:marco_mapa}
\end{figure}


\section{Fundamentos del aprendizaje profundo}
\label{sec:marco_fundamentos}

Esta sección presenta cómo una red neuronal produce una probabilidad por clase, cómo se ajustan sus
parámetros y cómo se comprueba que lo aprendido sirve fuera de los datos de entrenamiento. Solo incluye
lo que los capítulos siguientes usan sin volver a definirlo.

\subsection{Neurona artificial y perceptrón multicapa}
\label{sec:marco_perceptron}

Una neurona artificial calcula una suma ponderada de sus entradas, le suma un sesgo y aplica a ese
resultado una función no lineal, llamada activación. Una capa reúne varias neuronas que reciben la misma
entrada, y un perceptrón multicapa apila capas de modo que la salida de cada una es la entrada de la
siguiente. La capa $l$ transforma el vector $\mathbf{h}^{(l-1)}$ que recibe en
\begin{equation}
    \mathbf{h}^{(l)} = \phi\!\left(W^{(l)} \mathbf{h}^{(l-1)} + \mathbf{b}^{(l)}\right),
    \label{eq:marco_capa}
\end{equation}
donde $W^{(l)}$ es la matriz de pesos, $\mathbf{b}^{(l)}$ el vector de sesgos y $\phi$ la activación,
aplicada elemento a elemento. Sin la activación, la composición de capas se reduciría a una única
transformación lineal; con ella, la red representa funciones no lineales de la
entrada~\cite{rumelhart1986learning,goodfellow2016deep}.

Los pesos y los sesgos de todas las capas son los parámetros de la red, que se agrupan en $\theta$, y
entrenarla consiste en elegir sus valores. El perceptrón reaparece dentro de arquitecturas más complejas:
la cabeza de clasificación que convierte las características extraídas por una red visual en un puntaje
por clase es una capa de la forma de la Ecuación~\ref{eq:marco_capa}.

\subsection{Clasificación: softmax y entropía cruzada}
\label{sec:marco_softmax}

En un problema de $C$ clases, la última capa devuelve un vector de $C$ puntajes reales sin restricción,
los \emph{logits} $\mathbf{z}$. La función \emph{softmax} los convierte en una distribución de probabilidad,
\begin{equation}
    p_c = \frac{\exp(z_c)}{\sum_{j=1}^{C} \exp(z_j)}, \qquad c = 1, \dots, C,
    \label{eq:marco_softmax}
\end{equation}
con valores positivos que suman uno, y la clase predicha es la de mayor probabilidad. La pérdida de
entropía cruzada mide la probabilidad que el modelo asignó a la clase correcta $y$,
$\mathcal{L} = -\log p_y$. Vale cero cuando toda la probabilidad está en la clase correcta y crece sin
límite a medida que esa probabilidad se acerca a cero~\cite{goodfellow2016deep}.

Varias piezas del trabajo operan sobre estas cantidades. Las modificaciones de la pérdida para el
desbalance de clases (Sección~\ref{sec:marco_transferencia}) cambian la entropía cruzada o los logits
antes de la softmax. La calibración (Sección~\ref{sec:marco_eval_calib}) pregunta si las probabilidades
$p_c$ coinciden con las frecuencias de acierto observadas. Y una de las formas de obtener una categoría de
un modelo generativo, descrita en la Sección~\ref{sec:met_clasif_vlm}, compara las probabilidades que ese
modelo asigna al nombre de cada clase.

\subsection{Entrenamiento por descenso del gradiente}
\label{sec:marco_entrenamiento}

Entrenar la red es buscar los parámetros $\theta$ que minimizan la pérdida promedio sobre el conjunto de
entrenamiento. El descenso del gradiente lo hace de forma iterativa. Calcula el gradiente de la pérdida
respecto de los parámetros y los mueve un paso en la dirección opuesta,
\begin{equation}
    \theta \leftarrow \theta - \eta \, \nabla_{\theta} \mathcal{L},
    \label{eq:marco_gradiente}
\end{equation}
donde $\eta$ es la tasa de aprendizaje. En la variante estocástica, cada paso usa el gradiente estimado
sobre un lote de ejemplos en lugar del conjunto completo, y una época es una pasada por todos los
lotes~\cite{goodfellow2016deep}.

La retropropagación calcula ese gradiente para todos los parámetros a la vez
(Figura~\ref{fig:marco_perceptron}). La pasada hacia adelante evalúa la red capa por capa y guarda los
valores intermedios; la pasada hacia atrás aplica la regla de la cadena desde la pérdida hacia la entrada
y obtiene el gradiente de cada capa a partir del de la capa siguiente. Su costo es del mismo orden que el
de la pasada hacia adelante, y eso hace posible entrenar redes con millones de
parámetros~\cite{rumelhart1986learning,goodfellow2016deep}. A cambio, los valores intermedios y los
gradientes ocupan memoria durante el entrenamiento.

% Figura: perceptrón multicapa, pasada hacia adelante y retropropagación.
% Se usa en sec:marco_fundamentos, subsección de entrenamiento.

\begin{figure}[htbp]
    \centering
    \begin{tikzpicture}[x=1mm, y=1mm,
        neurona/.style = {circle, draw=figtrazo, minimum size=6mm, inner sep=0pt, font=\footnotesize},
    ]

        % Capas.
        \foreach \i in {1, 2, 3} {
            \node[neurona, fill=figazul!15, draw=figazul!80!black] (x\i) at (0, 16 - \i*8) {$x_{\i}$};
        }
        \foreach \i in {1, 2, 3, 4} {
            \node[neurona, fill=fignaranja!18, draw=fignaranja!70!black] (h\i) at (25, 20 - \i*8) {};
        }
        \foreach \i in {1, 2, 3} {
            \node[neurona, fill=figverde!12, draw=figverde!70!black] (z\i) at (50, 16 - \i*8) {$z_{\i}$};
        }

        % Conexiones.
        \foreach \i in {1, 2, 3} { \foreach \j in {1, 2, 3, 4} {
            \draw[figtrazo!60, thin] (x\i) -- (h\j);
        } }
        \foreach \j in {1, 2, 3, 4} { \foreach \k in {1, 2, 3} {
            \draw[figtrazo!60, thin] (h\j) -- (z\k);
        } }

        % Softmax y pérdida.
        \node[figtarea, text width=16mm, minimum height=8mm] (sm) at (72, 0) {softmax};
        \node[figsalida, text width=20mm, minimum height=8mm] (loss) at (100, 0)
            {pérdida\\$\mathcal{L} = -\log p_y$};
        \draw[figflecha] (56, 0) -- (sm.west);
        \draw[figflecha] (sm.east) -- node[figetiqueta, above] {$p_c$} (loss.west);

        % Rótulos de capas.
        \node[font=\footnotesize, text=figtrazo] at (0, -22) {entrada};
        \node[font=\footnotesize, text=figtrazo, align=center] at (25, -22) {capa oculta\\$\phi(W\mathbf{x}+\mathbf{b})$};
        \node[font=\footnotesize, text=figtrazo] at (50, -22) {\emph{logits}};

        % Pasadas.
        \draw[figflecha, figazul!80!black, thick] (0, 24) -- node[figetiqueta, above, text=figazul!80!black]
            {pasada hacia adelante: se calcula la salida} (100, 24);
        \draw[-{Latex[length=2mm, width=1.6mm]}, thick, densely dashed, figbermellon]
            (100, -32) -- node[figetiqueta, below, text=figbermellon]
            {retropropagación: se calcula $\nabla_\theta \mathcal{L}$ capa por capa} (0, -32);

    \end{tikzpicture}
    \caption[Perceptrón multicapa y retropropagación]{Perceptrón multicapa con una capa oculta. La pasada
    hacia adelante transforma la entrada capa por capa hasta los \emph{logits}, que la softmax convierte en
    probabilidades, y la pérdida mide la probabilidad asignada a la clase correcta. La retropropagación
    recorre el camino inverso con la regla de la cadena y obtiene el gradiente de la pérdida respecto de
    cada parámetro.}
    \label{fig:marco_perceptron}
\end{figure}


La tasa de aprendizaje única de la Ecuación~\ref{eq:marco_gradiente} trata igual a todos los
parámetros. Adam adapta el paso a cada uno: mantiene dos promedios móviles por parámetro, uno del
gradiente y otro de su cuadrado, y divide el primero por la raíz del segundo, de modo que los parámetros
con gradientes grandes o ruidosos avanzan con pasos más cortos~\cite{kingma2015adam}. Esos dos estados por
parámetro, sumados al gradiente, son el costo de memoria del ajuste fino que la
Sección~\ref{sec:marco_llm_lora} busca reducir.

\subsection{Generalización y regularización}
\label{sec:marco_regularizacion}

La pérdida de entrenamiento es un medio. Lo que interesa es el desempeño sobre casos que el modelo no
vio, y una red con suficientes parámetros puede memorizar el conjunto de entrenamiento: su pérdida sigue
bajando mientras el error sobre datos nuevos sube. Ese fenómeno se llama
sobreajuste~\cite{goodfellow2016deep}.

Para medirlo, los datos se dividen en tres particiones disjuntas. La de entrenamiento ajusta los
parámetros; la de validación elige entre configuraciones, hiperparámetros y épocas; y la de test se usa
una sola vez, al final, para estimar el desempeño sobre datos nuevos. Una decisión tomada mirando el test
lo convierte en parte del ajuste y vuelve optimista su estimación~\cite{goodfellow2016deep}. Por la misma
razón ningún caso puede aparecer en dos particiones, condición que HAM10000 no cumple en su partición
original (Sección~\ref{sec:met_particiones}).

Las técnicas de regularización reducen el sobreajuste restringiendo lo que la red puede aprender. El
decaimiento de pesos penaliza los parámetros grandes~\cite{goodfellow2016deep}. El \emph{dropout} apaga al
azar una fracción de las neuronas en cada paso de entrenamiento e impide que la red dependa de
combinaciones particulares de ellas~\cite{srivastava2014dropout}. El aumento de datos, que se detalla en
la Sección~\ref{sec:marco_transferencia}, amplía el conjunto con variaciones de las imágenes disponibles.
Conservar los parámetros de la época con mejor desempeño en validación, en lugar de los de la última,
cumple una función parecida~\cite{goodfellow2016deep}.

\subsection{Familias de redes según la estructura de la entrada}
\label{sec:marco_familias}

El perceptrón multicapa trata su entrada como un vector sin estructura: cada neurona se conecta con todas
las entradas y no distingue cuáles son vecinas. Las arquitecturas especializadas incorporan la estructura
de los datos. Las redes neuronales convolucionales (\gls{cnn}) aprovechan que en una imagen los píxeles
cercanos están relacionados y aplican los mismos filtros en toda la imagen~\cite{lecun1998gradient}; se
desarrollan en la Sección~\ref{sec:marco_cnn}. Las redes recurrentes procesan secuencias, como un texto,
un elemento por vez~\cite{elman1990finding}.

Una red recurrente mantiene un estado que resume lo visto hasta el momento y lo actualiza con cada nuevo
elemento, de modo que cada salida depende de toda la secuencia anterior~\cite{elman1990finding}. Para
entrenarla, el gradiente se propaga hacia atrás a lo largo de la secuencia, y en secuencias largas tiende
a desvanecerse, lo que dificulta aprender relaciones entre elementos lejanos. Las redes de memoria a corto
y largo plazo agregan compuertas que controlan qué se conserva en el estado y atenúan ese
problema~\cite{hochreiter1997long}. Aun así, el acceso a una posición lejana sigue pasando por todas las
intermedias, y el cálculo de un paso espera al del anterior, lo que impide paralelizar a lo largo de la
secuencia~\cite{vaswani2017attention}.

\section{Transformers y mecanismo de atención}
\label{sec:marco_atencion}

La atención reemplaza el recorrido paso a paso de las redes recurrentes por una comparación directa, en
la que cada posición de la salida pondera los valores de todas las posiciones de la entrada según la
compatibilidad entre una consulta y una clave. Como todas las posiciones se comparan a la vez, el cálculo
se paraleliza a lo largo de la secuencia. La arquitectura \emph{transformer} se construye sobre
esa operación, y de ella derivan tanto los modelos visuales de la Sección~\ref{sec:marco_vision} como
los modelos de lenguaje de la Sección~\ref{sec:marco_llm}~\cite{vaswani2017attention}.

La atención escalada por producto punto proyecta la entrada en tres espacios, consulta $Q$, clave
$K$ y valor $V$, y devuelve una combinación de valores pesada por la compatibilidad entre consulta
y clave:
\begin{equation}
    \mathrm{Attention}(Q, K, V) = \mathrm{softmax}\!\left(\frac{QK^{\top}}{\sqrt{d_k}}\right) V,
    \label{eq:marco_atencion}
\end{equation}
donde $d_k$ es la dimensión de la clave. La división por su raíz evita que el producto punto crezca
con la dimensión y sature la softmax, que dejaría de discriminar entre
posiciones~\cite{vaswani2017attention}. La Figura~\ref{fig:marco_atencion} muestra la matriz de pesos
que resulta para una secuencia corta.

% Figura: la matriz de pesos de atención.
% Se usa en sec:marco_atencion, después de la Ecuación de la atención escalada por producto punto.
% Los pesos son ilustrativos: muestran la forma de la operación, no un modelo entrenado.

\begin{figure}[htbp]
    \centering
    \begin{tikzpicture}[x=1mm, y=1mm]

        % Consultas (filas) y claves (columnas).
        \foreach \i in {1, 2, 3, 4} {
            \node[fignodo, minimum height=6mm, minimum width=9mm, inner sep=1pt,
                  draw=figazul!80!black, fill=figazul!12] (q\i) at (0, 30 - \i*8) {\footnotesize $q_{\i}$};
            \node[fignodo, minimum height=6mm, minimum width=9mm, inner sep=1pt,
                  draw=fignaranja!70!black, fill=fignaranja!18] (k\i) at (8 + \i*8, 34) {\footnotesize $k_{\i}$};
        }

        % Pesos softmax(QK^T / sqrt(d_k)); cada fila suma uno.
        \foreach \row/\vals in {1/{60,20,10,10}, 2/{15,55,20,10}, 3/{10,25,50,15}, 4/{5,10,25,60}} {
            \foreach \v [count=\col] in \vals {
                \fill[figceleste!\v!white] (4 + \col*8, 26 - \row*8) rectangle (12 + \col*8, 34 - \row*8);
                \node[font=\scriptsize] at (8 + \col*8, 30 - \row*8) {0.\ifnum\v<10 0\fi\v};
            }
        }
        \draw[figtrazo, shift={(12, -6)}] (0, 0) grid[step=8] (32, 32);

        % Salida.
        \foreach \i in {1, 2, 3, 4} {
            \node[fignodo, minimum height=6mm, minimum width=9mm, inner sep=1pt,
                  draw=figverde!70!black, fill=figverde!10, densely dashed] (o\i) at (66, 30 - \i*8)
                  {\footnotesize $o_{\i}$};
            \draw[figflecha] (45, 30 - \i*8) -- (o\i.west);
        }

        \node[font=\footnotesize, text=figtrazo] at (66, 32) {$o_i = \sum_j a_{ij} v_j$};
        \node[font=\footnotesize, text=figtrazo, align=center] at (28, -12)
            {pesos $a_{ij}$ = $\mathrm{softmax}\!\left(q_i k_j^{\top} / \sqrt{d_k}\right)$};

    \end{tikzpicture}
    \caption[Pesos de atención entre las posiciones de una secuencia]{Pesos de atención entre las
    posiciones de una secuencia de cuatro tokens. Cada fila compara la consulta $q_i$ con las claves de
    todas las posiciones, y la softmax convierte esas compatibilidades en pesos que suman uno; la salida
    $o_i$ es la combinación de los valores $v_j$ con esos pesos. Los valores de la matriz son ilustrativos.}
    \label{fig:marco_atencion}
\end{figure}


La atención multi-cabeza proyecta $Q$, $K$ y $V$ en varios subespacios de menor dimensión, aplica
la Ecuación~\ref{eq:marco_atencion} en cada uno por separado y concatena los resultados. Cada cabeza
aprende a capturar relaciones de distinto tipo entre posiciones, y el conjunto amplía la capacidad
de representación sin aumentar la dimensión total del modelo~\cite{vaswani2017attention}.

El transformer original, planteado para traducción automática, combina dos módulos. El codificador
contextualiza la entrada completa, y el decodificador genera la salida token a token atendiendo a lo
codificado y a lo ya generado. Las arquitecturas que usa este trabajo conservan uno solo de los dos:
el ViT de la Sección~\ref{sec:marco_vit} usa el codificador, y los modelos de lenguaje generativos de
la Sección~\ref{sec:marco_llm_grandes}, el
decodificador~\cite{vaswani2017attention,dosovitskiy2021image,brown2020language}.

\section{Aprendizaje profundo para visión}
\label{sec:marco_vision}

El primer componente del \emph{pipeline} recibe la imagen dermatoscópica y devuelve una probabilidad
para cada una de las siete categorías. Para esa tarea el trabajo enfrenta dos familias de arquitectura
de visión. Las redes convolucionales componen
filtros locales apilados y extraen una jerarquía de características. Los transformers visuales
tratan la imagen como una secuencia de tokens y la procesan con la atención de la
Sección~\ref{sec:marco_atencion}. DenseNet y ConvNeXt pertenecen a la primera familia y \gls{vit} y
Swin a la segunda, y las cuatro se comparan en el Capítulo~\ref{cp:res_clasif}.

\subsection{Redes convolucionales}
\label{sec:marco_cnn}

La operación de convolución desliza un filtro de pequeño tamaño sobre la entrada y calcula en cada
posición el producto punto entre el filtro y la ventana que cubre. La conectividad local y el uso de
un mismo filtro en toda la imagen son las dos decisiones de diseño del operador. La primera supone
que la información relevante se agrupa en regiones pequeñas, y la segunda, que un mismo patrón puede
aparecer en cualquier parte (Figura~\ref{fig:marco_convolucion}). Una CNN apila esas operaciones con
activaciones no lineales y reducciones de resolución intermedias~\cite{lecun1998gradient}.

% Figura: la operación de convolución.
% Se usa en sec:marco_cnn, después del primer párrafo.

\begin{figure}[htbp]
    \centering
    \begin{tikzpicture}[x=1mm, y=1mm]

        % Entrada 6x6 con la ventana de 3x3 resaltada.
        \fill[fignaranja!30] (5, 10) rectangle (20, 25);
        \draw[figtrazo!50] (0, 0) grid[step=5] (30, 30);
        \draw[figtrazo, thick] (0, 0) rectangle (30, 30);
        \draw[fignaranja!70!black, very thick] (5, 10) rectangle (20, 25);
        \node[font=\footnotesize, text=figtrazo] at (15, -5) {entrada $6\times 6$};

        % Filtro 3x3.
        \fill[fignaranja!30] (45, 7.5) rectangle (60, 22.5);
        \draw[fignaranja!70!black, shift={(45, 7.5)}] (0, 0) grid[step=5] (15, 15);
        \draw[fignaranja!70!black, very thick] (45, 7.5) rectangle (60, 22.5);
        \node[font=\footnotesize, text=figtrazo, align=center] at (52.5, -5) {filtro $3\times 3$\\(pesos compartidos)};
        \node[font=\large, text=figtrazo] at (37.5, 15) {$\ast$};

        % Salida 4x4 con la celda resultante.
        \fill[figverde!35] (80, 10) rectangle (85, 15);
        \draw[figtrazo!50] (75, 0) grid[step=5] (95, 20);
        \draw[figtrazo, thick] (75, 0) rectangle (95, 20);
        \draw[figverde!70!black, very thick] (80, 10) rectangle (85, 15);
        \node[font=\footnotesize, text=figtrazo, align=center] at (85, -5) {mapa de activación $4\times 4$};
        \node[font=\large, text=figtrazo] at (67.5, 15) {$=$};

        % De la ventana a la celda.
        \draw[figflecha, fignaranja!70!black, bend left=25] (20, 25) to
            node[figetiqueta, above] {producto punto} (80, 15);

    \end{tikzpicture}
    \caption[La operación de convolución]{La operación de convolución. El filtro recorre la entrada y, en
    cada posición, el producto punto entre sus pesos y la ventana que cubre produce una celda del mapa de
    activación. El mismo filtro se aplica en todas las posiciones, de modo que un patrón se detecta
    dondequiera que aparezca.}
    \label{fig:marco_convolucion}
\end{figure}


La profundidad compone esos rasgos locales en rasgos de mayor alcance: los bordes producen texturas,
las texturas producen partes de objetos y las partes producen objetos completos. VGG y ResNet
fijaron el repertorio de bloques que las arquitecturas posteriores refinan, con convoluciones
pequeñas apiladas en el primer caso y conexiones residuales en el segundo, y ese repertorio es el
punto de partida de las redes convolucionales
contemporáneas~\cite{simonyan2015very,he2015residual}.

En DenseNet cada capa recibe las características de todas las anteriores, además de las de la
inmediata, y aporta las suyas a todas las posteriores. Esa reutilización alcanza una capacidad
comparable con menos parámetros que las arquitecturas de conectividad simple, y con conjuntos del
orden de diez mil imágenes, como HAM10000, la economía de parámetros reduce el
sobreajuste~\cite{huang2018densely,he2015residual}.

ConvNeXt moderniza una ResNet adoptando decisiones de diseño de los transformers: convoluciones de
mayor tamaño en lugar de bloques pequeños apilados, menos no-linealidades y normalización por capa.
Con esas modificaciones, una red puramente convolucional iguala al Swin Transformer en exactitud y
escalabilidad. En el comparativo, ConvNeXt representa a las redes convolucionales de la misma época
que el ViT, de modo que el contraste entre familias enfrenta arquitecturas
contemporáneas~\cite{liu2022convnet}.

\subsection{Vision Transformer y Swin Transformer}
\label{sec:marco_vit}

El ViT divide la imagen en una rejilla de parches cuadrados, aplana cada parche y lo proyecta
a un token, y suma a cada token un \emph{embedding} de posición. Sobre esa secuencia aplica un
codificador transformer sin ninguna operación convolucional, y decide la categoría a partir de un
token de clase que se agrega al inicio de la secuencia~\cite{dosovitskiy2021image,vaswani2017attention}.

% Figura 2.1. Del parche al token en el Vision Transformer.
% Se usa en sec:marco_vision, subsección "Vision Transformer y Swin Transformer", después del
% párrafo que define el ViT.

\begin{figure}[htbp]
    \centering
    \begin{tikzpicture}[x=1mm, y=1mm]

        % Imagen dividida en una rejilla de parches.
        \draw[figtrazo, thick, fill=figrelleno] (-26, -12) rectangle (0, 12);
        \foreach \x in {-19.5, -13, -6.5} { \draw[figtrazo] (\x, -12) -- (\x, 12); }
        \foreach \y in {-6, 0, 6} { \draw[figtrazo] (-26, \y) -- (0, \y); }
        \node[font=\small, text=figtrazo, align=center] at (-13, -18) {Imagen dividida\\en parches};

        % Operación de proyección.
        \node[figtarea, text width=30mm] (proj) at (24, 0)
            {Proyectar cada parche\\a un token};

        % Secuencia de tokens.
        \foreach \i in {0, 1, 2, 3, 4} {
            \node[figsalida, minimum width=6mm, minimum height=6mm, inner sep=1pt]
                (t\i) at (52 + \i*9, 6) {\footnotesize $x_{\i}$};
            \node[font=\footnotesize, text=figtrazo] at (52 + \i*9, -3) {$+$};
            \node[font=\footnotesize, text=figtrazo] at (52 + \i*9, -9) {$p_{\i}$};
        }
        \node[font=\small, text=figtrazo, align=center] at (70, -18)
            {Secuencia de tokens con\\embedding de posición por parche};

        % Flechas.
        \draw[figflecha] (0, 0) -- (proj.west);
        \draw[figflecha] (proj.east) -- (46, 0) |- (t0.west);

    \end{tikzpicture}
    \caption[Del parche al token en el Vision Transformer]{Del parche al token en el Vision
    Transformer. La imagen se divide en una rejilla de parches cuadrados y cada parche se aplana y
    proyecta a un token. A cada token se le suma un embedding de posición, porque la operación de
    atención por sí sola no distingue el orden, y la secuencia entra al codificador
    \emph{transformer} sin volver a la retícula de píxeles.}
    \label{fig:marco_vit_patches}
\end{figure}


El ViT no incorpora las dos suposiciones que la convolución da por sentadas, la localidad y la
equivariancia a traslaciones, y tiene que aprenderlas de los datos. Pre-entrenado solo sobre
ImageNet, rinde por debajo de redes convolucionales de tamaño comparable; pre-entrenado sobre
corpus de entre 14 y 300 millones de imágenes, las supera. La variable que decide entre una familia
y otra es la escala del pre-entrenamiento, y por eso importa con qué pesos parte cada arquitectura
antes del ajuste sobre el conjunto médico~\cite{dosovitskiy2021image}.

El Swin Transformer atiende solo dentro de ventanas locales, con costo lineal en el número de
parches en lugar de cuadrático, y desplaza las ventanas entre capas para que la información cruce
de una a otra. Además reduce la resolución por etapas y devuelve una pirámide de mapas de
características, con la estructura jerárquica de una red convolucional y la operación de atención de
un transformer~\cite{liu2021swin}. La Figura~\ref{fig:marco_swin} muestra el desplazamiento de las
ventanas entre dos capas.

% Figura: ventanas y ventanas desplazadas del Swin Transformer. Adaptada de liu2021swin (Figura 2).
% Se usa en sec:marco_vit, en el párrafo del Swin.

\begin{figure}[htbp]
    \centering
    \begin{tikzpicture}[x=1mm, y=1mm]

        % Capa l: cuatro ventanas de 4x4 parches.
        \foreach \c/\x/\y in {figazul/0/16, fignaranja/16/16, figverde/0/0, figpurpura/16/0} {
            \fill[\c!22] (\x, \y) rectangle (\x + 16, \y + 16);
        }
        \draw[figtrazo!40] (0, 0) grid[step=4] (32, 32);
        \draw[figtrazo, very thick] (0, 0) rectangle (32, 32);
        \draw[figtrazo, very thick] (16, 0) -- (16, 32) (0, 16) -- (32, 16);
        \node[font=\footnotesize, text=figtrazo] at (16, -5) {capa $l$: ventanas regulares};

        % Capa l+1: ventanas desplazadas dos parches.
        \begin{scope}[xshift=50mm]
            \foreach \c/\x/\y/\w/\h in {
                figazul/0/24/8/8, fignaranja/8/24/16/8, figverde/24/24/8/8,
                figpurpura/0/8/8/16, figceleste/8/8/16/16, figbermellon/24/8/8/16,
                figazul/0/0/8/8, fignaranja/8/0/16/8, figverde/24/0/8/8} {
                \fill[\c!22] (\x, \y) rectangle (\x + \w, \y + \h);
            }
            \draw[figtrazo!40] (0, 0) grid[step=4] (32, 32);
            \draw[figtrazo, very thick] (0, 0) rectangle (32, 32);
            \draw[figtrazo, very thick] (8, 0) -- (8, 32) (24, 0) -- (24, 32) (0, 8) -- (32, 8) (0, 24) -- (32, 24);
            \node[font=\footnotesize, text=figtrazo] at (16, -5) {capa $l+1$: ventanas desplazadas};
        \end{scope}

        \draw[figflecha] (36, 16) -- (46, 16);

    \end{tikzpicture}
    \caption{Ventanas y ventanas desplazadas del Swin Transformer. Adaptada de \textcite{liu2021swin}.}
    \label{fig:marco_swin}
\end{figure}


\subsection{Transferencia de aprendizaje y desbalance de clases}
\label{sec:marco_transferencia}

Entrenar desde cero una de esas arquitecturas exige más imágenes de las que tiene un conjunto
médico típico. La transferencia de aprendizaje resuelve la escasez aprendiendo primero
representaciones sobre un corpus grande de imágenes naturales y reutilizándolas después para la
tarea de destino. La extracción de características congela el \emph{backbone} y solo entrena la
cabeza de clasificación; el ajuste fino actualiza también los pesos de la red, con una tasa de
aprendizaje menor~\cite{pan2010survey}.

Todas las arquitecturas del trabajo parten de pesos pre-entrenados sobre ImageNet, el corpus de
referencia del reconocimiento visual~\cite{russakovsky2015imagenet}. Con ese punto de partida común,
las diferencias que se miden entre ellas en el Capítulo~\ref{cp:res_clasif} se pueden atribuir a la
arquitectura.

El aumento de datos amplía el corpus de entrenamiento con transformaciones geométricas y
fotométricas de las imágenes disponibles: rotaciones, recortes, desplazamientos y cambios de color.
Funciona como regularizador y reduce el sobreajuste, pero, aplicado por igual a todas las clases,
deja la proporción entre ellas donde estaba~\cite{shorten2019survey}.

La distribución de clases de HAM10000 es muy asimétrica, y un modelo entrenado sin corrección
favorece a las clases mayoritarias~\cite{tschandl2018ham10000}. El Capítulo~\ref{cp:res_clasif}
compara tres modificaciones de la función de pérdida, que actúan durante el entrenamiento y dejan
intacta la métrica de evaluación: ponderar la pérdida por la inversa de la frecuencia de cada clase;
la pérdida focal, que atenúa el aporte de los ejemplos que el modelo ya clasifica con confianza; y
el \emph{balanced softmax}, que suma a cada logit el logaritmo de la frecuencia de su clase en el
entrenamiento~\cite{lin2017focal,ren2020balanced}.

Con esto el clasificador queda completo. Recibe la imagen y devuelve siete probabilidades, ninguna de
las cuales es un texto; el reporte exige un modelo que genere secuencias de palabras, que es el tema de
la sección siguiente.

\section{Modelos de lenguaje}
\label{sec:marco_llm}

Un modelo de lenguaje predice el siguiente token del texto a partir de los anteriores, y de esa
única tarea salen el pre-entrenamiento sobre corpus grandes, el ajuste a una tarea concreta y la
generación de textos nuevos. El generador del \emph{pipeline} pertenece a esta familia. Para el
trabajo interesan cómo se obtiene el modelo, cómo produce el texto y cuánto cuesta ajustarlo.

\subsection{Grandes modelos de lenguaje}
\label{sec:marco_llm_grandes}

Un \gls{llm} se obtiene por pre-entrenamiento auto-supervisado: sobre corpus de miles de
millones de palabras, el modelo aprende a predecir el siguiente token sin anotación humana. De esa
fase sale una distribución de probabilidad sobre el vocabulario condicionada a cualquier prefijo, y
esa distribución es la herramienta que después se reaprovecha para todo lo
demás~\cite{brown2020language,touvron2023llama}.

Dos familias se distinguen por el modo de usar el transformer. La de solo codificador entrena una
representación bidireccional del texto y sirve para clasificar y para comparar; la de solo
decodificador modela la secuencia en una dirección y sirve para generar. El componente de lenguaje
del trabajo pertenece a la segunda~\cite{devlin2019bert,brown2020language,touvron2023llama}.

El comportamiento del modelo se ajusta después en dos etapas. El ajuste fino supervisado lo entrena
sobre pares instrucción--respuesta, y la alineación con \gls{rlhf} lo entrena a preferir las
respuestas que un evaluador humano marca como mejores. Un modelo solo pre-entrenado continúa el
texto que recibe, mientras que uno alineado responde a instrucciones, y el generador del trabajo
recibe su tarea en forma de instrucción~\cite{ouyang2022training}.

La generación recorre la distribución de salida token a token. La búsqueda codiciosa y la búsqueda
por haz eligen las continuaciones más probables y producen textos planos y repetitivos. El muestreo
introduce azar en esa elección, y la temperatura, el \emph{top-k} y el muestreo de núcleo deciden
cuánto. El muestreo de núcleo recorta el vocabulario al subconjunto mínimo que acumula la masa de
probabilidad fijada, y con eso evita la degeneración textual~\cite{holtzman2020curious}. Como la
estrategia de decodificación cambia el texto generado y todas las métricas que se calculen sobre
él, se fija y se documenta como un hiperparámetro más del generador.

\subsection{Ajuste eficiente de parámetros: LoRA}
\label{sec:marco_llm_lora}

El ajuste fino completo de un modelo grande está limitado por la memoria más que por el tiempo de
cómputo. Con el optimizador Adam, cada parámetro entrenado arrastra su gradiente y dos estados
adicionales, y para un modelo de miles de millones de parámetros eso no cabe en la memoria de una
GPU de consumo. Congelar el modelo por completo y entrenar solo una cabeza nueva evita el problema,
pero limita lo que el ajuste puede cambiar~\cite{hu2021lora}.

\gls{lora} ocupa un punto intermedio entre las dos opciones. Sobre los pesos congelados $W$ entrena
un incremento de bajo rango,
\begin{equation}
    W' = W + BA, \qquad B \in \mathbb{R}^{d\times r},\quad A \in \mathbb{R}^{r\times k},\quad
    r \ll \min(d, k),
    \label{eq:marco_lora}
\end{equation}
y solo actualiza $A$ y $B$, que se inicializan de modo que el incremento valga cero al empezar
(Figura~\ref{fig:marco_lora}). El rango $r$ controla la capacidad de adaptación, y valores pequeños
bastan para llevar el modelo a un dominio concreto~\cite{hu2021lora}.

% Figura: reparametrización de LoRA. Adaptada de hu2021lora (Figura 1).
% Se usa en sec:marco_llm_lora, después de la Ecuación de LoRA.

\begin{figure}[htbp]
    \centering
    \begin{tikzpicture}[x=1mm, y=1mm]

        \node[figentrada, text width=10mm, minimum height=7mm] (x) at (0, 0) {$\mathbf{x}$};

        % Rama congelada.
        \node[fignodo, text width=36mm, minimum height=24mm, fill=figrelleno] (W) at (-24, 25)
            {pesos preentrenados\\$W \in \mathbb{R}^{d\times k}$\\[2pt]{\footnotesize congelados}};

        % Rama entrenable: A reduce a r, B vuelve a d.
        \draw[draw=fignaranja!70!black, fill=fignaranja!25, thick] (8, 16) -- (32, 16) -- (26, 22) -- (14, 22) -- cycle;
        \node[font=\footnotesize] at (20, 19) {$A \in \mathbb{R}^{r\times k}$};
        \draw[draw=fignaranja!70!black, fill=fignaranja!25, thick] (14, 28) -- (26, 28) -- (32, 34) -- (8, 34) -- cycle;
        \node[font=\footnotesize] at (20, 31) {$B \in \mathbb{R}^{d\times r}$};
        \node[font=\footnotesize, text=fignaranja!70!black, align=left, anchor=west] at (34, 25)
            {entrenables\\rango $r \ll \min(d,k)$};

        % Suma y salida.
        \node[circle, draw=figtrazo, inner sep=1pt, font=\small] (sum) at (0, 46) {$+$};
        \node[figsalida, text width=26mm, minimum height=7mm] (h) at (0, 58) {$\mathbf{h} = W\mathbf{x} + BA\mathbf{x}$};

        \draw[figflecha] (x.north) -- ++(0, 4) -| (W.south);
        \draw[figflecha] (x.north) -- ++(0, 4) -| (20, 16);
        \draw[figflecha] (20, 22) -- (20, 28);
        \draw[figflecha] (W.north) |- (sum.west);
        \draw[figflecha] (20, 34) |- (sum.east);
        \draw[figflecha] (sum.north) -- (h.south);

    \end{tikzpicture}
    \caption{Reparametrización de LoRA con el incremento de bajo rango $BA$. Adaptada de \textcite{hu2021lora}.}
    \label{fig:marco_lora}
\end{figure}


Los parámetros entrenables y sus estados del optimizador pasan a ser los de $A$ y $B$, una fracción
mínima de los de $W$, y con eso el ajuste vuelve a caber en el presupuesto de una GPU de consumo. Al
terminar, el incremento de la Ecuación~\ref{eq:marco_lora} se suma a $W$ y el modelo resultante se
ejecuta igual que el original, sin ramas ni latencia
adicional~\cite{hu2021lora,yilmaz2025resourceefficient}.

Los modelos generativos que el trabajo ajusta se adaptan con LoRA, y por eso esa parte de la
comparación se hace entera en hardware de consumo. El rango, los módulos congelados y la función de
pérdida del ajuste van en el Capítulo~\ref{cp:metodos}.

Un modelo de lenguaje, sin embargo, solo recibe texto. Para que redacte el reporte a partir de la imagen
hay que conectar su entrada con un codificador visual, y de eso se ocupan los modelos visión--lenguaje.

\section{Modelos visión--lenguaje}
\label{sec:marco_vlm}

\subsection{Alineación contrastiva: CLIP}
\pendiente{Se explica cómo CLIP proyecta imagen y texto a un espacio común, y cómo ese espacio permite la
clasificación \emph{zero-shot} por similitud y el \emph{linear probe}. Es la base de DermLIP.}

\subsection{Modelos multimodales generativos}
\label{sec:marco_vlm_generativos}
\pendiente{Se describe la proyección de características visuales al espacio de un LLM (LLaVA) y el
\emph{visual instruction tuning}, base de LLaVA-Med, Mini-VLM y SkinVL. Se aclara que estos modelos no
reutilizan el esquema codificador--decodificador del transformer original: los tokens visuales proyectados
entran como prefijo a un LLM de solo decodificador.}

\subsection{Generación de texto a partir de imágenes médicas}
\label{sec:marco_vlm_generacion}
\pendiente{Se presenta la tarea central del trabajo: producir texto clínico a partir de una imagen. Se
describen el esquema \emph{image captioning} aplicado al dominio médico y su evolución hacia la redacción de
informes (radiología primero, dermatología después), qué datos requiere (pares imagen--texto clínico), y los
dos problemas característicos de estos modelos: la alucinación de hallazgos y el \emph{hedging}, es decir, la
tendencia a redactar descripciones genéricas y sin compromiso diagnóstico. También se distingue entre la
fidelidad diagnóstica del reporte y su calidad como texto, que retoma las dos condiciones de la
Figura~\ref{fig:intro_dos_tareas} (categoría acertada y descripción fiel a la imagen); sobre esa distinción
se apoya toda la evaluación del Capítulo~\ref{cp:res_gen}.}

\subsection{Generación aumentada por recuperación}
\pendiente{Se explica RAG en general y su variante basada en vecinos visuales (kNN sobre
embeddings).}

\section{Evaluación de modelos}
\label{sec:marco_evaluacion}

\subsection{Métricas de clasificación}
\label{sec:marco_eval_clasif}
\pendiente{Se definen exactitud, F1 ponderado y macro, sensibilidad macro, AUC, y sensibilidad y
especificidad para la tarea binaria, y se muestra por qué la exactitud es engañosa con clases
desbalanceadas.}

\subsection{Calibración de probabilidades}
\label{sec:marco_eval_calib}
\pendiente{Se presentan el diagrama de confiabilidad, ECE, MCE y Brier, y los métodos de calibración post-hoc
\emph{temperature scaling} y \emph{vector scaling}.}

\subsection{Métricas de generación de texto}
\label{sec:marco_eval_texto}
\pendiente{Se describen las métricas con referencia (BLEU, ROUGE-L, METEOR, CIDEr, BERTScore): qué
comparan entre el texto generado y el de referencia, y qué miden y qué no.}

Las métricas con referencia suponen que el texto de referencia es la descripción correcta de la imagen,
y en dermatología ese supuesto es débil. Cuando varios dermatólogos marcan por separado los mismos
criterios dermatoscópicos sobre las mismas imágenes, la concordancia va de pobre a
moderada~\cite{carrera2016validity}. El texto de referencia de un par imagen--texto recoge entonces una
lectura de la lesión entre varias posibles, y una descripción generada puede apartarse de él sin estar
equivocada. Una métrica que compara contra ese único texto la penaliza igual que a una descripción
errónea.

Las métricas sin referencia se apoyan en otra propiedad del dominio. La descripción de una lesión se
escribe con un vocabulario cerrado de términos para color, patrón, forma y distribución, que la
comunidad dermatoscópica estandarizó por consenso. Dentro de ese vocabulario, algoritmos como la lista
de 7 puntos convierten la presencia de ciertos hallazgos en peso diagnóstico, de modo que el camino
entre la imagen y la etiqueta pasa por entidades con
nombre~\cite{kittler2016standardization,argenziano1998checklist}. Con un inventario finito de
entidades se puede contar cuántas de ellas menciona un texto generado sin compararlo con otro texto.

\pendiente{Se definen las dos métricas sin referencia que usa el trabajo, la cobertura de entidades
clínicas y la consistencia diagnóstica, y qué mide y qué no mide cada familia.}

\subsection{Inferencia estadística para comparar modelos}
\pendiente{Se presentan los intervalos por \emph{bootstrap} pareado, la prueba de McNemar y la prueba de
Wilcoxon, y qué pregunta responde cada una. Para Wilcoxon se cita, junto al original, demsar2006statistical,
que justifica su uso para comparar clasificadores.}

\subsection{Interpretabilidad visual}
\pendiente{Se describen los mapas de activación (Grad-CAM) y sus limitaciones.}

\section{Sesgo y generalización}
\label{sec:marco_generalizacion}

\pendiente{Se define el cambio de distribución (\emph{distribution shift}): un modelo aprende la
distribución de su conjunto de entrenamiento, y su desempeño medido sobre una partición de ese mismo
conjunto no garantiza el desempeño sobre imágenes de otra procedencia. Falta la cita del concepto
general (propuesta en refs\_pendientes.md).}

En dermatología el cambio de distribución más visible es el de modalidad. La fotografía clínica
registra la lesión tal como se ve a simple vista, con la piel que la rodea, y la dermatoscopía la
examina con aumento e iluminación controladas, con lo que muestra estructuras que la fotografía no
resuelve, como la red pigmentaria y los patrones
vasculares~\cite{dinnes2018dermoscopy,kittler2016standardization}. Los rasgos visuales tienen
distribuciones distintas en uno y otro registro, y un modelo entrenado sobre una modalidad no se
traslada a la otra sin volver a medirlo. HAM10000 es dermatoscópico y SkinCaRe es de fotografía
clínica, de modo que toda comparación entre los dos conjuntos arrastra esa diferencia de
dominio~\cite{tschandl2018ham10000,shen2025skincare}.

\pendiente{Se presentan las disparidades de desempeño por tono de piel en los conjuntos públicos y la
cuestión de la generalización a poblaciones distintas de la de entrenamiento, que motiva la evaluación
en HIBA (Sección~\ref{sec:met_datasets}).}

\section{Estado del arte}
\label{sec:estado_arte}
\pendiente{Esta sección revisa los trabajos directamente relacionados con la tesina: diagnóstico automático
de lesiones de piel, modelos visión--lenguaje en dermatología y generación de reportes dermatológicos. Cada
bloque cierra indicando qué toma este trabajo de esa línea o en qué se diferencia.}

\subsection{Diagnóstico automático de lesiones de piel}
\pendiente{Un cuadro compara los conjuntos públicos del área (ISIC, PH2, Fitzpatrick17k, SkinCon, Derm1M y
los que usa el trabajo) en modalidad, tamaño, etiquetas, disponibilidad de texto y población; la descripción
detallada de los conjuntos usados está en la Sección~\ref{sec:met_datasets}. Se recorre la evolución desde las CNNs clásicas hacia los transformers, con los resultados de
referencia sobre ISIC y HAM10000 con los que se compara el Capítulo~\ref{cp:res_clasif}, y los problemas
recurrentes del área: desbalance, generalización y sesgo por fototipo.}

\subsection{Modelos visión--lenguaje en dermatología}
\pendiente{Se presentan los modelos pre-entrenados a gran escala con texto clínico (PanDerm, MONET, Derm1M y
DermLIP, MAKE) y su desempeño \emph{zero-shot}, revisión que justifica la elección de DermLIP. Se incluyen
también los VLMs médicos generales efectivamente evaluados en este trabajo o comparables con él (LLaVA-Med y
Phi-3.5-Vision) y la evidencia disponible sobre su desempeño en dermatología.}

\subsection{Generación de reportes dermatológicos}
\pendiente{Se revisan SkinGPT-4, SkinVL (MM-Skin), DermatoLlama/SCALEMED, la destilación de conocimiento de
Chakrabarti et al., Skin-R1, SkinGEN y los ensambles con LLM, junto con la forma en que cada uno evalúa
(métricas con referencia, evaluación por especialistas) y el hardware que requiere. Una tabla comparativa
resume modelo, tarea, datos, métricas y hardware.

La sección cierra con la brecha que aborda esta tesina, que es triple: faltan comparaciones bajo un mismo
protocolo entre clasificadores supervisados, VLMs generales y modelos visión--lenguaje dermatológicos; la
mayoría de los generadores se evalúa por similitud textual y no por fidelidad diagnóstica; y pocos trabajos
consideran la incertidumbre del diagnóstico o el hardware accesible.}

\chapter[Materiales y métodos]{Materiales y métodos}
\label{cp:metodos}
% Extensión estimada: 20--25 páginas.
% Fuentes: sección "Configuración experimental" y subsecciones "Diseño del experimento" del documento de
% experimentación, Propuesta (datasets, técnicas), planes en Codigo/docs/ideas/ (en particular
% split_lesion_semillas_plan_20260912.md). Incorpora el protocolo unificado pedido en el feedback del tutor
% del 12/09/2026 (Codigo/docs/feedback_20260912.md).
%
% Feedback del tutor (14/09/2026):
%  - "Particiones y auditoría de fuga de información" se renombra: el término no se entiende fuera de
%    contexto. Ahora la subsección se llama por lo que produce (la partición disjunta por lesión) y el
%    término "fuga de información (data leakage)" se introduce dentro del texto, definido.
%  - La subsección "Métricas" repetía el marco teórico. Se fusiona con la de selección: el marco teórico
%    define qué es cada métrica, aquí se justifica cuál se usa para decidir.
%  - "Robustez, calibración e interpretabilidad" no decía a qué se refería; se renombra por extenso.
%  - "Sistema propuesto" pasa a llamarse "pipeline", que es el término corriente en el área para la
%    arquitectura completa: clasificador, kNN y generador encadenados.
%  - Se incorpora aquí "Uso de un VLM generativo como clasificador", que estaba en el marco teórico pese a
%    ser parte del enfoque propio.
%
% Revisión 25/09/2026: recibe del marco teórico la prosa ya escrita sobre las categorías de HAM10000
% (§HAM10000), los costos asimétricos de confusión (§selección de métricas) y las partes del reporte
% (§generación de reportes). Ver fichas/02-marco-teorico.md, decisión 14.

\pendiente{Este capítulo describe todo lo necesario para reproducir los experimentos: datos, particiones,
modelos, \emph{pipeline} y protocolo de evaluación. Los resultados se presentan en los capítulos
siguientes.}

\section{Metodología de trabajo}
\label{sec:met_metodologia}
\pendiente{Se describe el ciclo iterativo \emph{hipótesis $\rightarrow$ ejecución $\rightarrow$ análisis
$\rightarrow$ decisión} con registro sistemático: un plan con hipótesis y criterio de éxito, la ejecución y
un resultado con veredicto binario. Varias etapas del trabajo nacieron de refutar la hipótesis anterior, de
modo que el orden de los capítulos de resultados sigue el orden en que se tomaron las decisiones: baseline
supervisado, VLMs en clasificación, refutación de la hipótesis del desbalance, especialización de dominio,
generación de reportes y ensamble en cascada, y análisis de rigor y robustez.}

\section{Conjuntos de datos y particiones}
\label{sec:met_datasets}

\subsection{HAM10000}
\label{sec:met_ham}

El clasificador asigna cada imagen a una etiqueta de un inventario cerrado, y ese inventario es el de
HAM10000, con siete categorías: \gls{mel}, \gls{nv}, \gls{bcc}, \gls{akiec}, \gls{bkl}, \gls{df} y
\gls{vasc}. Según su gravedad, NV, BKL, DF y VASC son benignas, AKIEC reúne lesiones premalignas y
carcinomas in situ, y MEL y BCC son malignas~\cite{tschandl2018ham10000}. Esa gradación determina
cuánto cuesta cada confusión del clasificador, y con ello qué métricas se usan para juzgarlo
(Sección~\ref{sec:met_seleccion}).

\pendiente{Se presentan la composición del conjunto, la distribución de clases y su desbalance.}

\subsection{Partición disjunta por lesión}
\label{sec:met_particiones}
\pendiente{HAM10000 contiene varias imágenes de una misma lesión y la partición original se estratificó por
imagen, de modo que imágenes de una misma lesión quedaron repartidas entre entrenamiento y test. Esta
subsección explica ese problema, conocido como fuga de información entre particiones (\emph{data leakage}),
lo cuantifica (el 36\% del test comparte lesión con entrenamiento; el 63\% en melanoma) y describe cómo se
midió su efecto sobre las métricas: reevaluación sobre el subconjunto de test cuya lesión no aparece en
entrenamiento ni en validación, y control con un modelo que nunca vio HAM10000 (DermLIP \emph{zero-shot}).

Se describe luego la partición disjunta por lesión que se construyó como reemplazo, con test idéntico al
original y validación agrupada (\texttt{StratifiedGroupKFold}), sus tamaños finales y qué resultados se
recalcularon sobre ella.}

\subsection{SkinCaRe, MM-Skin y HIBA}
\pendiente{Se describe el rol de cada conjunto: entrenamiento de la generación (SkinCAP y SkinCoT),
evaluación con referencia (MM-Skin) y evaluación fuera de distribución (HIBA, solo test). Se indica la
modalidad de cada uno, en particular que SkinCaRe es de fotografía clínica y HAM10000 dermatoscópico, con
remisión a la Sección~\ref{sec:marco_generalizacion}.}

\subsection{Preprocesamiento y aumento de datos}
\pendiente{Se detallan la resolución, la normalización, los aumentos de entrenamiento y el TTA.}

\section{Modelos evaluados}
\label{sec:met_modelos}

\subsection{Clasificadores supervisados}
\pendiente{Se describen DenseNet121, ViT-B/16, Swin-Tiny y ConvNeXt-Tiny: pre-entrenamiento, cabeza de
clasificación, pérdidas y estrategias de desbalance evaluadas.}

\subsection{Modelos visión--lenguaje médicos generales}
\pendiente{Se presentan Mini-VLM (Phi-3.5-Vision) y LLaVA-Med: arquitectura, tamaño y motivo de su
elección.}

\subsection{Modelo visión--lenguaje dermatológico: DermLIP}
\pendiente{Se describen la arquitectura y el pre-entrenamiento de DermLIP (Derm1M) y sus modos de uso:
\emph{zero-shot}, \emph{linear probe}, kNN y SVM sobre \emph{embeddings}.}

\subsection{Generador dermatológico: SkinVL}
\pendiente{Se presenta SkinVL y el motivo de incluirlo como generador especializado de comparación.}

\section{Clasificación con modelos visión--lenguaje}
\label{sec:met_clasif_vlm}
\pendiente{Obtener una clase de un modelo generativo admite varias formas, y la elección condiciona el
resultado. Esta sección presenta las cuatro estrategias evaluadas —respuesta libre (\emph{zero-shot}),
\emph{likelihood scoring} sobre las etiquetas candidatas, razonamiento estructurado (\emph{Chain-of-Thought})
y ajuste fino con LoRA— junto con los prompts, el parseo de respuestas y la configuración de cada una.}

\section{Generación de reportes}
\label{sec:met_generacion}

El reporte dermatológico deja por escrito la lectura de la lesión, y de sus partes el trabajo se ocupa
de dos: la descripción de lo que se ve y la impresión diagnóstica. Son las que corresponden a las dos
tareas del problema, y las que los conjuntos de pares imagen--texto del área anotan junto a cada
imagen~\cite{shen2025skincare}. El diagnóstico diferencial y la recomendación de manejo, que un reporte
clínico completo también lleva, quedan fuera del alcance.

\pendiente{Se describen el ajuste fino del generador sobre SkinCaRe con LoRA, el formato concreto del
reporte generado y la configuración de decodificación.}

\section{Pipeline propuesto: ensamble en cascada consciente de la incertidumbre}
\label{sec:met_pipeline}

\subsection{Cascada clasificador--generador (\emph{grounding})}
\pendiente{Se explica por qué separar el diagnóstico de la redacción, cómo se inyecta el diagnóstico en el
prompt y qué variantes se evalúan (clasificador, DermLIP y \emph{oracle} como cota superior). Un diagrama
resume el pipeline.}

\subsection{Grounding por recuperación}
\pendiente{Se describen la búsqueda de vecinos visuales con DermLIP (kNN, $k=3$) y la construcción del
contexto.}

\subsection{Ruteo por incertidumbre en tres niveles}
\pendiente{Se explica cómo el acuerdo entre clasificador y kNN determina el nivel de cada caso (afirmativo,
diferencial o abstención con derivación) y se presentan las plantillas de cada nivel.}

\section{Protocolo de evaluación}
\label{sec:met_protocolo}

\subsection{Selección de modelos, métricas y criterio de decisión}
\label{sec:met_seleccion}

Las confusiones entre categorías no tienen el mismo costo clínico. Confundir MEL con NV retrasa la
extirpación de una lesión cuya letalidad depende del estadio en que se diagnostica; confundir NV con MEL
lleva a extirpar una lesión benigna sin necesidad; las confusiones entre las clases benignas tienen menor
consecuencia~\cite{tschandl2018ham10000,nci2026melanoma}. Esa asimetría obliga a medir el desempeño clase
por clase, además de en conjunto, y a dar a las clases minoritarias el mismo peso que a la mayoritaria.

\pendiente{El Capítulo~\ref{cp:marco} define qué mide cada métrica; esta subsección indica cuáles se usan
para decidir y por qué. Se justifican la sensibilidad macro como métrica principal de clasificación y el F1
macro como complemento, el umbral de especificidad $\geq 0.85$ fijado en validación para la detección de
melanoma, las métricas de calibración y las dos de generación (Nivel 1 con referencia; Nivel 2 con
\emph{clinical entity recall} y \emph{diagnostic consistency}), con la distinción entre fidelidad diagnóstica
y validez clínica.

Se describen a continuación los barridos de hiperparámetros y el criterio de selección, que opera
enteramente en validación: el test no interviene en ninguna decisión. Hay que distinguir dos niveles. El
\emph{checkpoint} de cada corrida es la época con mejor F1 ponderado en validación. La corrida ganadora
entre las del barrido, en cambio, se elige por \emph{F1 macro} en validación, porque el F1 ponderado está
dominado por la clase mayoritaria ---NV es el 67\,\% de las imágenes--- y la sensibilidad macro, que sería
la alternativa natural por ser la métrica principal de este trabajo, ignora los falsos positivos y por lo
tanto premia a los modelos que sobre-predicen las clases minoritarias. El F1 macro es el único de los tres
que pondera las clases por igual y a la vez penaliza los falsos positivos. La sensibilidad macro se reporta
como resultado principal, pero no se usa para elegir.

Los resultados se reportan como media $\pm$ desvío entre semillas.}

\subsection{Protocolos de inferencia}
\pendiente{Una tabla única indica, para cada resultado reportado, si se usa vista única o TTA, qué
\emph{checkpoint} y qué partición.}

\subsection{Análisis estadístico}
\pendiente{Se indica qué intervalo o prueba corresponde a cada métrica y comparación: \emph{bootstrap}
pareado para diferencias de métricas agregadas, McNemar para aciertos por imagen y Wilcoxon para
\emph{recall} por clase.}

\subsection{Análisis complementarios: validación externa, perturbaciones, calibración e interpretabilidad}
\pendiente{Además de las métricas sobre el test, se realizan cuatro análisis que responden preguntas
distintas de ``cuánto acierta el modelo'': si el desempeño se sostiene en otra población (evaluación en
HIBA), si se sostiene bajo peores condiciones de captura (perturbaciones de adquisición: ruido, desenfoque,
compresión JPEG y pérdida de resolución), si las probabilidades que entrega el modelo son confiables
(calibración post-hoc ajustada en validación) y si el modelo mira la lesión o el contexto de la imagen
(Grad-CAM). Esta subsección describe el procedimiento de cada uno.}

\section{Infraestructura y reproducibilidad}
\label{sec:met_infra}
\pendiente{Se describen los recursos de cómputo (GPU local RTX 4070, clúster TOKO con RTX 4080 y Mendieta con
A30), el uso de SLURM, el seguimiento de experimentos y la organización del código y del registro de
ejecuciones. El detalle está en el Apéndice~\ref{ap:reproducibilidad}.}

\chapter[Resultados: clasificación y detección de melanoma]{Resultados: clasificación diagnóstica y detección de melanoma}
\label{cp:res_clasif}
% Extensión estimada: 20--25 páginas.
% Fuentes: sección "Clasificación diagnóstica" del documento de experimentación, con los números del split
% disjunto por lesión (Codigo/outputs/results/lesion_split_report.json, cerrado el 2026-09-15 con las
% semillas 42/43/44).
%
% Feedback del tutor (14/09/2026): la sección de cierre se llama "Discusión" en lugar de "Síntesis del
% capítulo". La discusión transversal, que cruza este capítulo con el de generación, queda en el
% Capítulo~\ref{cp:discusion}.

\pendiente{Este capítulo responde la PI1 y contrasta las hipótesis H1, H3 y H4. Los resultados se presentan
en el orden en que se tomaron las decisiones: cada sección presenta un resultado, lo interpreta y justifica
el paso siguiente.}

\section{Clasificadores supervisados}
\label{sec:res_supervisados}
\pendiente{Se presentan el barrido y el modelo elegido en validación, sus resultados en test como media
$\pm$ desvío entre semillas (exactitud, F1 ponderado y macro, sensibilidad macro y \emph{recall} por clase),
el efecto de las estrategias de desbalance y el efecto de corregir la fuga respecto del split original. Se
discute también la dispersión entre semillas, muy desigual según la métrica: la exactitud varía poco mientras
que el \emph{recall} de melanoma es la métrica más inestable, lo que obliga a reportar media y desvío y no
una corrida única.}

\section{Modelos visión--lenguaje médicos generales}
\label{sec:res_vlm_generales}
\pendiente{Se comparan las cuatro estrategias con Mini-VLM y LLaVA-Med frente al baseline supervisado, y se
muestra el colapso de las predicciones hacia una o dos clases.}

\section{Re-diagnóstico del colapso}
\label{sec:res_rediagnostico}
\pendiente{Se contrasta la hipótesis del desbalance con la evidencia que la refuta: el colapso aparece en
\emph{zero-shot}, los datos de ajuste no están sesgados hacia las clases problemáticas y cada configuración
colapsa hacia clases distintas.}

\section{La especialización de dominio recupera la discriminación}
\label{sec:res_dermlip}
\pendiente{Se presentan los resultados de DermLIP en \emph{zero-shot}, \emph{linear probe}, kNN y SVM (split
por lesión) frente a los VLMs generales y a los clasificadores supervisados.}

\section{Comparación con el estado del arte}
\label{sec:res_sota_clasif}
\pendiente{Se comparan los resultados con los publicados sobre HAM10000 (Sección~\ref{sec:estado_arte}),
señalando las diferencias de protocolo (particiones, fuga, métricas) que limitan la comparación.}

\section{Calibración y significancia estadística}
\label{sec:res_calib_stats}
\pendiente{Se analiza la calibración antes y después de \emph{temperature} y \emph{vector scaling}, junto con
el efecto de los pesos por clase, y se reportan los intervalos de confianza y las pruebas estadísticas entre
los modelos finales.}

\section{Robustez}
\label{sec:res_robustez}
\pendiente{Se evalúa la generalización a HIBA (otra población y otros fototipos) y a perturbaciones de las
condiciones de adquisición.}

\section{Interpretabilidad}
\label{sec:res_interpretabilidad}
\pendiente{Se analizan mapas Grad-CAM para determinar si el modelo atiende a la lesión.}

\section{Detección de melanoma}
\label{sec:res_melanoma}
\pendiente{Se evalúan las tareas binarias melanoma-vs-resto y maligno-vs-benigno (AUC, sensibilidad y
especificidad con umbral fijado en validación) por modelo y para el ensamble, y se comparan con referencias
clínicas y de la literatura.}

\section{Discusión}
\label{sec:res_clasif_discusion}
\pendiente{Se resume qué hipótesis se sostienen y cuáles no, y por qué el clasificador especializado pasa a
ser la fuente de diagnóstico del sistema de generación.}

\chapter[Resultados: generación de reportes]{Resultados: generación de reportes y pipeline integrado}
\label{cp:res_gen}
% Extensión estimada: 25--30 páginas.
% Fuentes: sección "Generación de reportes" del documento de experimentación,
% Codigo/docs/resul/incertidumbre_llava_replica_resul_20260912.md y la sección "Costo computacional" del
% mismo documento.
%
% Feedback del tutor (14/09/2026): el capítulo de costo computacional era demasiado corto para sostenerse
% solo (3--5 páginas) y se incorpora como última sección de medición de este capítulo, que es donde se mide
% el pipeline completo. La comparación de hardware contra el estado del arte, que es interpretación y no
% medición, pasa al Capítulo~\ref{cp:discusion}. La sección de cierre se llama "Discusión".

\pendiente{Este capítulo responde la PI2 y evalúa el OE2, junto con la parte del OE3 que corresponde a la
generación de reportes (calidad del texto, tratamiento de la incertidumbre y costo computacional). Los reportes se generaron
sobre el test original, que es idéntico al del split por lesión; la robustez frente a la fuga de información
se reporta sobre el subconjunto limpio.}

\section{Generación directa con VLMs ajustados}
\label{sec:res_gen_directa}
\pendiente{Se presentan los resultados de Nivel 1 y Nivel 2 de los generadores ajustados sin
\emph{grounding}, que producen descripciones ricas pero diagnósticamente poco fieles.}

\section{Reporte frente a razonamiento}
\label{sec:res_gen_razonamiento}
\pendiente{Se describe el hallazgo metodológico sobre el entrenamiento con narrativas de razonamiento
(SkinCoT) y el \emph{trade-off} de estilo que distorsiona las métricas con referencia.}

\section{Ensamble en cascada}
\label{sec:res_gen_cascada}
\pendiente{Se justifica el uso de un ensamble y se presentan sus resultados con distintas fuentes de
diagnóstico (clasificador, DermLIP y \emph{oracle} como cota superior), con dos generadores y por clase.}

\section{Taxonomía de errores}
\label{sec:res_gen_errores}
\pendiente{Se clasifican los errores de los reportes (diagnóstico incorrecto, insuficiente, repetición, otros)
y se analiza cómo cambian con cada arquitectura.}

\section{Un generador dermatológico: SkinVL}
\label{sec:res_gen_skinvl}
\pendiente{Se evalúa si un generador especializado en dermatología evita la necesidad de la cascada.}

\section{Aporte de la imagen y del VLM}
\label{sec:res_gen_ablaciones}
\pendiente{Las ablaciones sin imagen o con imagen cambiada muestran que la imagen no cambia qué diagnóstico se
redacta, pero sí qué se describe. Se analiza qué aporta el VLM frente a una plantilla (riqueza clínica,
especificidad por caso, lenguaje): el clasificador aporta el diagnóstico y el VLM, la descripción.}

\section{Grounding por recuperación}
\label{sec:res_gen_rag}
\pendiente{Se compara el contexto construido con vecinos visuales frente al \emph{grounding} por
clasificador.}

\section{Ensamble consciente de la incertidumbre}
\label{sec:res_gen_incertidumbre}
\pendiente{Se presentan la distribución de casos por nivel (afirmativo, diferencial, abstención) y la
fidelidad diagnóstica sobre los casos no derivados y sobre la cohorte completa, con intervalos, junto con la
réplica con LLaVA-Med. Se explica cómo leer el valor de consistencia, que no es una exactitud multiclase
porque incluye el diferencial y excluye los casos derivados. Se reporta la robustez con el clasificador y el
kNN del split por lesión y sobre el subconjunto limpio.}

\section{Ejemplos cualitativos}
\label{sec:res_gen_cualitativos}
\pendiente{Se muestran reportes seleccionados que ilustran aciertos, diferenciales, derivaciones y errores
típicos.}

% Sección condicionada a que se concrete la evaluación con dermatólogos (septiembre--octubre 2026). Si no se
% realiza, se elimina y queda como limitación.
% \section{Evaluación por especialistas}
% \label{sec:res_gen_especialistas}

\section{Costo computacional del pipeline}
\label{sec:res_costo}
\pendiente{Se reportan el tiempo y la memoria de entrenamiento o ajuste de cada modelo en las GPUs
utilizadas, y el tiempo y la memoria por imagen de cada componente del \emph{pipeline} (clasificador, kNN y
generador) junto con el costo por reporte del pipeline completo. Estos números cierran el eje transversal de
hardware accesible planteado en la introducción; su comparación con los requerimientos de los trabajos del
estado del arte se retoma en el Capítulo~\ref{cp:discusion}.}

\section{Discusión}
\label{sec:res_gen_discusion}
\pendiente{Se resume qué arquitectura resulta más fiel diagnósticamente y bajo qué condiciones.}

\chapter[Discusión]{Discusión general}
\label{cp:discusion}
% Extensión estimada: 8--12 páginas.
% Fuentes: sección "Discusión y conclusiones" del documento de experimentación (incluye "Amenazas a la
% validez" y "Limitaciones") y la comparación de hardware del antiguo capítulo de costo computacional.
%
% Feedback del tutor (14/09/2026): cada capítulo de resultados cierra con su propia discusión; este capítulo
% queda para lo transversal, lo que ninguno de los dos podría cerrar por sí solo. Por eso se fusionan aquí
% "Comparación con trabajos relacionados" dentro de las secciones temáticas y "Amenazas a la validez" con
% "Limitaciones".

\pendiente{Este capítulo interpreta en conjunto los resultados de los dos capítulos anteriores y los ubica
frente a la literatura, sin repetir sus tablas.}

\section{Respuesta a las preguntas de investigación}
\label{sec:disc_preguntas}
\pendiente{Se responden la PI1 y la PI2 con la evidencia de los capítulos de resultados y se resume el estado
final de cada hipótesis.}

\section{Especialización de dominio frente a modalidad}
\label{sec:disc_especializacion}
\pendiente{Se discute por qué el factor determinante no es que el modelo sea visión--lenguaje, sino que esté
especializado en dermatología, y cómo se ubica esa conclusión respecto de los trabajos revisados en la
Sección~\ref{sec:estado_arte}.}

\section{División de trabajo entre clasificador y generador}
\label{sec:disc_division}
\pendiente{Se discute que el clasificador aporta el diagnóstico y el VLM la descripción, y qué implica esto
para el diseño de sistemas de generación de reportes médicos.}

\section{Fidelidad diagnóstica y validez clínica}
\label{sec:disc_validez_clinica}
\pendiente{Se delimita qué establecen y qué no las métricas automáticas, y el papel de la incertidumbre y la
derivación en un sistema de apoyo.}

\section{Lecciones metodológicas}
\label{sec:disc_lecciones}
\pendiente{Se discuten la fuga de información por lesión en HAM10000 y su efecto sobre los resultados
publicados en el área, la importancia de seleccionar exclusivamente en validación, la inestabilidad entre
semillas de las métricas de las clases minoritarias y el sesgo de las métricas con referencia al comparar
estilos de reporte.}

\section{Viabilidad en hardware accesible}
\label{sec:disc_hardware}
\pendiente{Se relacionan el desempeño alcanzado y el costo medido en la Sección~\ref{sec:res_costo} con los
requerimientos de cómputo de los trabajos del estado del arte, para discutir qué parte de la brecha con el
SOTA se explica por escala de modelo y de datos y qué parte no.}

\section{Limitaciones y amenazas a la validez}
\label{sec:disc_limitaciones}
\pendiente{Se analizan las amenazas a la validez interna, externa, de constructo y de conclusión estadística,
junto con las limitaciones del trabajo: escala de los modelos, costo, ausencia de reportes redactados por
especialistas, generalización fuera de distribución, anclajes y generadores no evaluados, y generación de
reportes sobre el test original.

Para cada hipótesis exploratoria (H3 a H6) se indica si hubo evidencia independiente de los datos que la
motivaron (HIBA, la partición disjunta por lesión, semillas nuevas) o si se contrastó sobre el mismo
material, en cuyo caso su respaldo queda como exploratorio. La Sección~\ref{sec:intro_preguntas} remite
aquí.}

\chapter[Conclusiones]{Conclusiones}
\label{cp:conclusiones}
% Extensión estimada: 4--6 páginas. Sin información nueva.
% Feedback del tutor (14/09/2026): al eliminarse la sección "Contribuciones" de la introducción, los aportes
% del trabajo se enuncian acá, ya respaldados por los resultados.

\pendiente{Este capítulo cierra el trabajo con sus conclusiones y el grado de cumplimiento de los objetivos.}

\section{Conclusiones generales}
\label{sec:concl_generales}
\pendiente{Se sintetizan los hallazgos principales en pocos párrafos.}

\section{Cumplimiento de los objetivos}
\label{sec:concl_objetivos}
\pendiente{Una tabla relaciona cada objetivo específico con el capítulo que lo aborda y el resultado
obtenido.}

\section{Aportes del trabajo}
\label{sec:concl_aportes}
\pendiente{Se enuncian los aportes, cada uno con remisión al capítulo que lo sustenta: una comparación
sistemática bajo un mismo protocolo entre clasificadores supervisados, VLMs médicos generales y un modelo
visión--lenguaje dermatológico; el diagnóstico del colapso de los VLMs generales en clasificación
dermatológica y la evidencia de que el factor determinante es la especialización de dominio y no la
modalidad; un ensamble en cascada consciente de la incertidumbre que mejora la fidelidad diagnóstica de los
reportes, replicado con dos generadores; la auditoría de fuga de información entre particiones de HAM10000 y
los resultados corregidos con una partición disjunta por lesión; y el análisis de viabilidad sobre hardware
de consumo.}


%%% Bibliografía %%%
\renewcommand{\refname}{Bibliografía}
\printbibliography[title={\refname},heading=bibintoc]

%%% Apéndices %%%
\appendix
\ifthenelse{\equal{\LanguageOption}{spanish}}{
    \addtocontents{toc}{\protect\contentsline{chapter}{Apéndices}{}{}}
}{
    \addtocontents{toc}{\protect\contentsline{chapter}{Appendices}{}{}}
}

\ifthenelse{\equal{\MediaOption}{paper}}{\blankpage}{\clearpage}
\begin{center}
    \crimsonfont
    \thispagestyle{empty}
        
    \vspace*{\fill}
    \ifthenelse{\equal{\LanguageOption}{spanish}}{%
        {\LARGE\fontsize{26}{26}\selectfont\textcolor{maincolor}{Apéndices}\par}
    }{%
        {\LARGE\fontsize{26}{26}\selectfont\textcolor{maincolor}{Appendices}\par}
    }
    \vspace*{\fill}
\end{center}
\MediaOptionLogicBlank
\chapter{Hiperparámetros y configuración de entrenamiento}
\label{ap:hiperparametros}
\pendiente{Este apéndice detalla las grillas de los barridos de Swin-Tiny y ConvNeXt-Tiny, los
hiperparámetros de DenseNet121 y ViT-B/16, las configuraciones de las pérdidas frente al desbalance y la
configuración de LoRA de los VLMs ajustados.}

\chapter{Prompts y plantillas}
\label{ap:prompts}
\pendiente{Este apéndice reúne los prompts de clasificación con VLMs (\emph{zero-shot}, \emph{likelihood},
CoT), los prompts de generación, las plantillas de \emph{grounding} para los niveles afirmativo, diferencial
y abstención, y el diccionario de entidades clínicas usado en las métricas de Nivel 2.}

\chapter{Resultados complementarios}
\label{ap:complementarios}
% Material disponible: lesion_split_report.json (ranking de 58+ trials, 14 corridas de desbalance/legacy,
% semillas), leakage_analysis.json (por clase), P03_perclass.png, P01_calibration.png,
% P24_robustness_examples.png, dermlip_split_eval.json.

\pendiente{Este apéndice presenta resultados detallados que complementan los
Capítulos~\ref{cp:res_clasif} y~\ref{cp:res_gen}:
\begin{itemize}
  \item el ranking completo de los trials del barrido en validación, con su resultado en test;
  \item los resultados por semilla de los modelos elegidos;
  \item la tabla completa de estrategias frente al desbalance para Swin-Tiny, ConvNeXt-Tiny y DenseNet121;
  \item las métricas por clase y matrices de confusión de todos los clasificadores;
  \item la auditoría de fuga de información desagregada por clase;
  \item los diagramas de confiabilidad y los ejemplos de perturbaciones de adquisición.
\end{itemize}}

% \chapter{Rúbrica de evaluación clínica}
\label{ap:rubrica}
% Fuente: Codigo/docs/rubrica_validacion_clinica.md. No se incluye en tesina.tex hasta que se concrete la
% evaluación con especialistas.
\pendiente{Este apéndice presenta la rúbrica usada para la evaluación de reportes por especialistas, los
criterios de selección de la muestra y las instrucciones para los evaluadores.}
 % Pendiente: se incorpora si se concreta la evaluación con especialistas.
\chapter{Reproducibilidad}
\label{ap:reproducibilidad}
\pendiente{Este apéndice describe la estructura del repositorio, el entorno de software, los scripts
principales y los comandos para reproducir cada tabla, así como el uso de los clústeres (envío de trabajos
y sincronización de resultados).}

\section{Semillas}
\label{sec:ap_semillas}

Todos los entrenamientos, la partición por lesión y la generación de los reportes fijan la semilla 42. Esa es también
la semilla de la ejecución de referencia de cada clasificador supervisado (Sección~\ref{sec:met_estadistica}), y las
otras dos ejecuciones (Sección~\ref{sec:met_supervisados_ensamble}) usan las semillas 43 y 44. Como la generación usa
decodificación voraz, la semilla no cambia los reportes, y se fija solo para que el entorno sea el mismo en todas las
ejecuciones.

\chapter{Fichas de conjuntos de datos y modelos externos}
\label{ap:fichas}
% Era el único anexo; se unificó con los apéndices (revisión del tutor, 2026-10-05).
\pendiente{Este apéndice reúne las fichas descriptivas y licencias de los conjuntos de datos utilizados
(HAM10000, SkinCaRe, MM-Skin, HIBA y Derm1M) y las tarjetas de los modelos preentrenados empleados
(DermLIP, LLaVA-Med, Phi-3.5-Vision y SkinVL), tal como las publican sus autores.}


% Anexos de la plantilla UMU: se unificaron con los apéndices (revisión del tutor, 2026-10-05).
% \ifthenelse{\equal{\LanguageOption}{spanish}}{
    \addtocontents{toc}{\protect\contentsline{chapter}{Anexos}{}{}}
}{
    \addtocontents{toc}{\protect\contentsline{chapter}{Annexes}{}{}}
}

\setcounter{chapter}{11} % To start at the "L" chapter.
\MediaOptionLogicAnnexes
\begin{center}
    \crimsonfont
    \thispagestyle{empty}
        
    \vspace*{\fill}
    \ifthenelse{\equal{\LanguageOption}{spanish}}{%
        {\LARGE\fontsize{26}{26}\selectfont\textcolor{maincolor}{Anexos}\par}
    }{%
        {\LARGE\fontsize{26}{26}\selectfont\textcolor{maincolor}{Annexes}\par}
    }
    \vspace*{\fill}
\end{center}
\MediaOptionLogicBlank

%%% Contraportada %%%
\ifthenelse{\equal{\MediaOption}{paper}}{\blankpage}{}

\clearpage
\null
\thispagestyle{empty}

% Back page: blanca, sin el fondo de la plantilla UMU, para acompañar la carátula.

\newgeometry{margin=1.98cm, top=1.47cm, bottom=1.47cm}
\noindent\clearpage
\restoregeometry

\end{document}
